% !TeX spellcheck = fr_FR
\documentclass[10pt,a5paper,twoside]{book}
\usepackage[utf8]{inputenc}
\usepackage[french]{babel}
\usepackage[T1]{fontenc}
% \usepackage{multicol} % pour {\'e}crire dans certaines zones en colonnes : \begin{multicols}{nb colonnes}...\end{multicols}
\usepackage{multirow}
\usepackage{geometry}
\geometry{
    a5paper,
    top=1.5cm,
    bottom=1.5cm,
    left=1.5cm,
    right=1.5cm,
    headsep=0.5cm,
    footskip=1cm
}

%  Packages pour la mise en page et le style 
\usepackage{parskip}
\usepackage{setspace}
\onehalfspacing
\usepackage{lipsum}
\usepackage{fancyhdr}
\pagestyle{fancy}
\fancyhf{}
\fancyhead[LE,RO]{\thepage}
\fancyhead[RE,LO]{\nouppercase{\leftmark}}
\renewcommand{\headrulewidth}{0.4pt}
\renewcommand{\footrulewidth}{0pt}

%  Packages pour les images et les couleurs 
\usepackage{graphicx}
\usepackage{xcolor}
\definecolor{darkred}{RGB}{128,0,0}
\definecolor{darkblue}{RGB}{0,0,128}

%  Packages pour les tableaux 
\usepackage{array}
\usepackage{multirow}
\usepackage{booktabs}
\newcolumntype{P}[1]{>{\raggedright\arraybackslash}p{#1}}

%  Packages pour les encadrés 
\usepackage{framed}
\usepackage{color}
\definecolor{shadecolor}{RGB}{240,240,240}

%  Packages pour les listes 
\usepackage{enumitem}
\setlist[itemize]{leftmargin=*}

%  Packages pour les hyperliens 
\usepackage{hyperref}
\hypersetup{
    colorlinks=true,
    linkcolor=darkblue,
    urlcolor=darkred,
    citecolor=darkblue
}

%  Titre et auteur 
\title{\textbf{Livret des Agences Occultes \\ France et Europe \\ pour Delta Green}}
\author{}
\date{}

\begin{document}

\maketitle

\newpage

\tableofcontents

\newpage

%  Introduction 
\chapter*{Introduction}
\addcontentsline{toc}{chapter}{Introduction}

Ce livret est conçu pour être utilisé avec les règles de \textbf{Delta Green}, afin d'ajouter des agences françaises et européennes spécialisées dans la lutte contre l'occulte. Vous y trouverez des descriptions détaillées des agences, des fiches de personnages, des idées de scénarios, ainsi que des localisations et descriptions pour enrichir vos parties.

\begin{shaded}
	\begin{center}
		\textbf{Note :} Ce document est inspiré des travaux de Tristan Lhomme, notamment de l'agence Félicie, et intègre des éléments de l'univers de Delta Green. Les règles de base de Delta Green sont nécessaires pour jouer. \\ \textbf{Crédits :} Gabriel Chandesris, inspiré par Tristan Lhomme, Delta Green (Arc Dream Publishing).
	\end{center}
\end{shaded}

\newpage

%  Agence Félicie 
\chapter{L'Agence Félicie}

\section{Description Générale}

L'\textbf{Agence Félicie} est une organisation secrète française fondée en \textbf{1945}, spécialisée dans la lutte contre les menaces surnaturelles et occultes. Successeur du \textbf{Bureau S} (actif pendant la Seconde Guerre mondiale), elle opère sous la supervision directe du \textbf{ministère de la Défense} ou de la \textbf{DGSE}. Officiellement dissoute dans les \textbf{années 1980}, elle continue d'exister clandestinement, avec des agents infiltrés dans divers services de l'état.

\begin{shaded}
	\textbf{Devises :} <<Protéger la France, même de l'ineffable. >> \\ <<Le secret est notre bouclier, la discrétion notre arme. >>
\end{shaded}

\subsection{Contexte Historique}

Félicie a été créée à la fin de la Seconde Guerre mondiale pour enquêter sur les phénomènes paranormaux et les activités occultes qui menacent la sécurité nationale. Elle a hérité des archives et des ressources du Bureau S, une organisation similaire active pendant la guerre. Son existence est \textbf{niée par l'état français}, et ses agents sont souvent déclarés \textbf{morts au combat} ou \textbf{démissionnaires} en cas de disparition.

\subsection{Niveau de Classification}
\begin{shaded}
	\textbf{Niveau de Secret :} \textbf{Défense Nationale -- Accès Réservé (Omega-1)}
	\begin{itemize}
		\item \textbf{Accès restreint} : Seuls les agents triés sur le volet ont connaissance de son existence.
		\item \textbf{Compartimentage strict} : Chaque agent n'a accès qu'à une partie des informations, selon son niveau d'habilitation.
		\item \textbf{Effacement des traces} : Les missions sont niées par l'état en cas de fuite. Les agents disparus sont déclarés morts au combat ou démissionnaires.
		\item \textbf{Protocole de désinformation} : Utilisation de couvertures (fausses identités, missions officielles bidon).
	\end{itemize}
\end{shaded}

\subsection{Structure Organisationnelle}

Félicie est organisée en \textbf{cellules autonomes}, sans hiérarchie centrale visible, pour limiter les risques de trahison ou d'infiltration.

\begin{center}
	\begin{tabular}{|>{\bfseries}P{4cm}|P{7cm}|}
	\hline
	\textbf{Division}							&	\textbf{Rôle}	\\ \hline
	Direction Générale (Conseil des Anciens)	&	Supervise les opérations, prend les décisions stratégiques. Composée d'anciens militaires, hauts fonctionnaires, et scientifiques. \\ \hline
	Bureau des Opérations						&	Planifie et coordonne les missions sur le terrain. \\ \hline
	Service Recherche et Analyse				&	Étudie les phénomènes occultes, analyse les artefacts, veille stratégique. \\ \hline
	Service Contre-Espionnage Occulte			&	Détecte et neutralise les infiltrations liées à des organisations occultes. \\ \hline
	Service Technologie et Équipements			&	Développe et gère les outils spécifiques (armes, artefacts, technologies occultes). \\ \hline
	Service Nettoyage							&	Efface les preuves, élimine les témoins, couvre les scandales. \\ \hline
	\end{tabular}
\end{center}

\newpage

\subsection{Protocoles de Sécurité}

%% \vfill

\begin{itemize}
	\item \textbf{Désinformation} : Félicie utilise des couvertures pour masquer ses activités (ex : missions de "sécurité culturelle" pour justifier des déplacements).
	\item \textbf{Neutralisation} : Les agents trop exposés ou corrompus sont \textbf{éliminés} (suicides maquillés, accidents).
	\item \textbf{Contrôle mental} : Certains agents subissent des conditionnements psychologiques pour éviter les fuites (hypnose, drogues amnésiques).
	\item \textbf{Protocole Janus} : En cas de compromission extrême, activation d'un \textbf{bombardement ciblé} ou d'une \textbf{opération de désinformation massive}.
\end{itemize}

\dotfill \vfill \dotfill

\subsection{Profil des Agents}
Les agents de Félicie sont recrutés parmi des profils variés, souvent après avoir été exposés à des phénomènes inexplicables. Le recrutement se fait \textbf{par la force ou par nécessité} : on leur propose de rejoindre l'agence ou de "disparaître".

\begin{shaded}
	\textbf{Exemple de recrutement :}
	Un agent de la DGSE, témoin d'un phénomène inexplicable en mission, est approché par Félicie. On lui propose de "rejoindre l'équipe" ou de "devenir un regret de la République". La plupart choisissent de servir, par patriotisme ou par peur.
\end{shaded}

\newpage

\begin{center}
\rotatebox{90}{
	\begin{tabular}{|>{\bfseries}P{3cm}|>{\bfseries}P{3cm}|P{5cm}|P{3cm}|}
	\hline
	\textbf{Type}			&	\textbf{Origine}				&	\textbf{Compétences}											&	\textbf{Faiblesses} \\ \hline
	Le Baroudeur			&	Légion étrangère, RAID, DGSE	&	Combat, survie, infiltration, résistance à la folie				&	Paranoïa, alcoolisme, syndrome de stress post-traumatique \\ \hline
	Le Scientifique			&	CNRS, CEA, universités			&	Recherche, analyse d'artefacts, protocoles occultes				&	Obsession pour le savoir interdit, risque de corruption \\ \hline
	Le Policier				&	Police nationale, gendarmerie	&	Enquête, interrogatoire, gestion de crise						&	Cynisme, désillusion, burnout \\ \hline
	L'Occultiste Repenti	&	Ancien membre de culte, sorcier	&	Connaissance des rituels, détection des menaces surnaturelles	&	Instabilité mentale, tentations de replonger dans le mal \\ \hline
	Le Diplomate			&	Quai d'Orsay, ambassades		&	Négociation, couverture internationale							&	Manipulateur, risque de trahison \\ \hline
	Le Nettoyeur			&	Services secrets, mercenaires	&	élimination discrète, falsification de preuves					&	Froid, déshumanisé \\ \hline
	\end{tabular} }
\end{center}

\newpage

\subsection{Méthodes et Outils}
Félicie utilise un mélange de \textbf{technologies modernes} et de \textbf{protocoles occultes} pour lutter contre les menaces.

\begin{itemize}
	\item \textbf{Armes} :
	\begin{itemize}
		\item Pistolets silencieux (ex : \textbf{Heckler \& Koch MK 23}).
		\item Armes bénites (balles en argent, lames en fer froid).
		\item Grenades à fragmentation "spéciales" (efficaces contre les entités).
	\end{itemize}
	\item \textbf{Équipement} :
	\begin{itemize}
		\item Détecteurs de magie (basés sur des artefacts anciens).
		\item Caméras thermiques (pour voir les entités).
		\item Kits de survie en milieu hostile.
	\end{itemize}
	\item \textbf{Artefacts} :
	\begin{itemize}
		\item Grimoires protégés (conservés dans des coffres scellés).
		\item Objets maudits sous surveillance (utilisés en dernier recours).
		\item Talismans de protection (ex : \textbf{Médaille de Saint Michel}, efficace contre les démons).
	\end{itemize}
	\item \textbf{Protocoles} :
	\begin{itemize}
		\item Rituels de confinement (pour neutraliser des entités).
		\item Exorcisme militaire (adapté des textes religieux anciens).
		\item Protocole "Silence Radio" : En cas de contact avec une entité, interdiction totale de communication non sécurisée.
	\end{itemize}
	\item \textbf{Réseaux} :
	\begin{itemize}
		\item Collaboration avec des sociétés secrètes (franc-maçonnerie, Rosicruciens).
		\item Contacts dans les \textbf{archives nationales} et les \textbf{musées} (ex : Louvre, Bibliothèque nationale).
		\end{itemize}
	\item \textbf{Bases secrètes} :
	\begin{itemize}
		\item Sous-sols du \textbf{ministère de la Défense} (Paris).
		\item Châteaux isolés (ex : \textbf{Château de Puymartin} en Dordogne).
		\item Bunkers abandonnés (ex : \textbf{Ligne Maginot} en Alsace).
	\end{itemize}
\end{itemize}

\newpage

\subsection{Menaces Typiques}
Félicie affronte des menaces uniques, souvent liées à l'histoire et au folklore français.

\begin{itemize}
\item \textbf{Cultes} :
	\begin{itemize}
		\item \textbf{Les Héritiers de Baphomet} : Culte médiéval ressuscité, actif dans le \textbf{Languedoc}.
		\item \textbf{La Confrérie du Graal Noir} : Société secrète cherchant à corrompre les artefacts chrétiens.
		\item \textbf{Les Enfants de Chaugnar Faugn} : Culte lovecraftien infiltré dans les \textbf{universités françaises}.
	\end{itemize}
	\item \textbf{Entités} :
	\begin{itemize}
		\item \textbf{Marie-Morgane} : Sirène maléfique liée à la légende de la \textbf{ville d'Ys} (Bretagne).
		\item \textbf{Korrigans} : Créatures malicieuses des \textbf{forêts bretonnes}.
		\item \textbf{L'Ankou} : Personnification de la Mort en \textbf{Bretagne}, capable de manipuler les ombres.
		\item \textbf{Les Lavandières de la Nuit} : Fantômes de blanchisseuses, actives près des \textbf{rivières}.
	\end{itemize}
	\item \textbf{Phénomènes paranormaux} :
	\begin{itemize}
		\item \textbf{Châteaux hantés} (ex : \textbf{Château de Brissac}, "le géant vert").
		\item \textbf{Forêts enchantées} (ex : \textbf{Forêt de Brocéliande}, liée aux légendes arthuriennes).
		\item \textbf{Objets maudits} (ex : \textbf{Le Miroir de la Duchesse de Montpensier}, conservée au \textbf{Musée de Cluny}).
	\end{itemize}
\end{itemize}

\newpage

\subsection{Scénarios}

\subsubsection{Ethor bedi zure erresuma (La Saga Félicie)}

\textbf{Contexte :} Années 1970, \textbf{Laos/Cambodge/Pays Basque}.
\begin{itemize}
	\item \textbf{Acte 1} : Les agents de Félicie partent enquêter sur un \textbf{trafic d'objets d'art maudits} entre le Laos et le Cambodge. Ils infiltrent un réseau de contrebandiers.
	\item \textbf{Acte 2} : L'enquête les mène à un \textbf{temple en pleine jungle}, où un culte tente de réveiller un \textbf{Grand Ancien}. Ils doivent éviter les militaires \textbf{Viêt-Congs} et les pièges du culte.
	\item \textbf{Acte 3} : Les agents sont \textbf{capturés} par des sectateurs et se réveillent \textbf{30 ans plus tard} au Pays Basque, avec un \textbf{trou de mémoire}. Ils sont accusés de trahison par leur propre organisation.
	\item \textbf{Dilemme final} :
	\begin{itemize}
		\item Détruire le culte et \textbf{mourir} en héros.
		\item Laisser le Grand Ancien se réveiller et \textbf{condamner l'humanité}.
		\item Utiliser un \textbf{rituel interdit} pour prolonger leur vie (au prix de leur âme).
	\end{itemize}
\end{itemize}

\begin{shaded}
\textbf{Note du MJ :} Ce scénario est conçu pour être \textbf{psychologiquement intense}. Les joueurs doivent gérer leur \textbf{perte de mémoire}, leur \textbf{culpabilité} (ils ont peut-être commis des atrocités pendant leurs 30 ans de "disparition"), et leur \textbf{loyauté envers Félicie}.
\end{shaded}

\newpage

\subsubsection{Le Crépuscule du Bureau S}

\textbf{Contexte :} France, \textbf{années 1940} (Seconde Guerre mondiale).
\begin{itemize}
	\item Les joueurs incarnent des agents du \textbf{Bureau S}, précurseur de Félicie, en pleine \textbf{déliquescence}.
	\item \textbf{Mission} : Enquêter sur des \textbf{phénomènes paranormaux} liés à l'occupation nazie (ex : \textbf{expériences occultes} menées par la \textbf{Ahnenerbe}).
	\item \textbf{Dilemme} : Sauver la France ou \textbf{laisser les nazis activer un rituel} qui pourrait les détruire eux-mêmes ?
\end{itemize}

\subsubsection{Aktion Hel}
\textbf{Contexte :} \textbf{Alsace, 1940}, ligne Maginot.
\begin{itemize}
	\item Un \textbf{meurtre impossible} a lieu dans une casemate de la ligne Maginot.
	\item Les agents doivent découvrir que l'assassin est un \textbf{fantôme} ou une \textbf{créature surnaturelle}.
	\item \textbf{Ambiance} : Huis clos, paranoïa, et horreur à la \textbf{Alien}.
\end{itemize}

\newpage


%  Autres possibilités en France 
\chapter{Autres idées pour une Agence française}

\section{Cadre institutionnel}

En France, l'agence pourrait s'intégrer comme une branche secrète au sein du ministère de l'Intérieur ou de la Défense, ou comme une organisation autonome sous la supervision directe du Premier ministre, à l'instar de la DGSE ou de la Communauté française du renseignement (CFR). Son financement serait assuré par l'État, avec un budget dédié et des ressources partagées avec les services de renseignement existants.

L'agence bénéficierait d'un cadre légal strict, encadré par des lois sur le secret d'État et la sécurité / sûreté nationale, lui donnant un niveau d'autorité élevé mais limité par des contrôles parlementaires et judiciaires. Elle pourrait collaborer avec des institutions scientifiques comme le CEA ou l'Ifri pour la recherche sur les phénomènes paranormaux.

\newpage

\section{Propositions}

\begin{center}
	\begin{tabular}{|>{\bfseries}P{4cm}|P{7cm}|}
	\hline
	\textbf{Nom proposé }													&	\emph{Justification historique, culturelle ou linguistique}	\\ \hline
	Clos Barbisier															&	Jardin public à Besançon, nommé d'après un personnage folklorique local, évoquant un lieu mystérieux et secret.	\\ \hline
	Office Central de Lutte Contre le Trafic des Biens Culturels (OCBC)		&	Service de police judiciaire français, évoquant la lutte contre des crimes liés à des biens culturels, souvent associés à des phénomènes occultes.	\\ \hline
	Groupe Indépendant d'Études Ésotériques de Paris						&	Association fondée en 1890 par Gérard Encausse (Papus), dédiée à l'étude de l'ésotérisme, incarnant une tradition occulte française.	\\ \hline
	Mansur al-Amir Bi-Ahkamillah											&	10e calife fatimide, connu pour sa lutte contre des sectes occultes, évoquant une figure historique liée à la lutte contre l'occulte.	\\ \hline
	Bureau des Affaires Non Conventionnelles (BANC)							&	Nom fictif évoquant la gestion d'affaires secrètes et paranormales, avec une sonorité administrative française.	\\ \hline
	BELPHÉGOR																&	Bureau d'Expertise Libre sur les Phénomènes Hautement Esoteriques sur les Grandes Œuvres Retrouvées.	\\ \hline
	\end{tabular}
\end{center}

\newpage

\section{Organisation interne}

L'agence serait structurée en plusieurs divisions spécialisées :
\begin{itemize}
	\item \textbf{Direction générale} : supervisant l'ensemble des opérations, coordonnant avec les autres services de renseignement.  
	\item \textbf{Division Recherche et Analyse} : chargée de l'étude des phénomènes occultes, de la collecte de données et de la veille stratégique.  
	\item \textbf{Division Intervention} : composée d'agents de terrain, spécialisés dans la neutralisation des menaces surnaturelles.  
	\item \textbf{Division Contre-espionnage occulte} : dédiée à la détection et à la neutralisation des infiltrations ou manipulations liées à des organisations occultes.  
	\item \textbf{Division Technologie et Équipements} : développant et gérant les outils spécifiques, artefacts et technologies adaptées à la lutte contre l'occulte.
\end{itemize}

Le recrutement se ferait parmi des profils variés : anciens militaires, policiers, scientifiques, universitaires spécialisés en sciences humaines ou exactes, ainsi que des experts en sécurité informatique et en investigation criminelle.

\section{Profil des agents}

Les agents seraient issus de milieux divers : forces armées, police, gendarmerie, services de renseignement, mais aussi du monde académique et scientifique. Ils devraient posséder des compétences pointues dans leur domaine, mais aussi une grande résilience psychologique, car ils seraient confrontés à des horreurs indicibles.

Les profils psychologiques incluraient une forte capacité d'analyse, un esprit critique, une grande discrétion, mais aussi une ouverture d'esprit face à des phénomènes inexpliqués. Certains agents pourraient avoir des liens avec des traditions ésotériques françaises (martinisme, hermétisme), ce qui pourrait être un atout dans la compréhension des menaces occultes.

\section{Méthodes et outils}

L'agence utiliserait un mélange de technologies modernes (surveillance électronique, analyse de données, cyberdéfense) et de protocoles occultes (rituels défensifs, utilisation d'artefacts anciens). Elle s'appuierait sur des experts civils pour des consultations spécifiques, et sur des bases de données secrètes regroupant des connaissances sur les phénomènes paranormaux.

Les agents seraient formés à des techniques d'enquête rigoureuses, à la gestion de crises impliquant des entités surnaturelles, et à la protection de la population contre des menaces invisibles.

\section{Menaces typiques}

Les menaces incluraient des cultes liés à l'histoire médiévale française, des sociétés secrètes inspirées de la franc-maçonnerie occulte, des entités surnaturelles issues du folklore local (dahut, korrigans, lavandières, ankou), ainsi que des phénomènes paranormaux liés à des lieux historiques (châteaux hantés, forêts enchantées).

\section{Défis uniques}

La méfiance historique envers les sociétés secrètes en France, la complexité du cadre légal et la nécessité de maintenir le secret tout en assurant la sécurité nationale seraient des défis majeurs. La coordination entre les différents services de renseignement et la gestion des informations sensibles serait également cruciale.

\section{Localisations pour scénarios}

ÀParis : MNHN (Muséeum National d'Histoire Naturelle... et Surnaturelle !), Louvre, Musée de l'homme, Musée du quai Branly... 

%  Agence Européenne EORA 
\chapter{L'Agence Européenne EORA}

\section{Description Générale}

L'\textbf{EORA (European Occult Response Agency)} est une organisation autonome créée par l'\textbf{Union Européenne} pour lutter contre les menaces occultes à l'échelle continentale. Elle opère sous l'autorité du \textbf{Centre de Renseignement de l'UE (IntCen)} ou d'\textbf{Europol}.

\begin{shaded}
\textbf{Devises :} "Unis contre l'ineffable." \\ "La connaissance est notre arme, le secret notre bouclier."
\end{shaded}

\subsection{Contexte Historique}
L'EORA a été fondée pour répondre aux menaces transnationales qui dépassent les capacités des agences nationales (comme Félicie). Elle collabore avec les services de renseignement des états membres, tout en maintenant une certaine autonomie. Son existence est \textbf{niée par les institutions européennes}, et ses opérations sont financées par un \textbf{budget secret}.

\subsection{Niveau de Classification}
\begin{shaded}
\textbf{Niveau de Secret :} \textbf{UE Classifié -- Niveau Omega-1}
	\begin{itemize}
		\item \textbf{Accès} : Seuls les états membres peuvent désigner des agents.
		\item \textbf{Compartimentage} : Chaque pays ne connaît qu'une partie des opérations.
		\item \textbf{Financement} : Budget UE Classifié, avec contributions nationales.
		\item \textbf{Juridiction} : L'EORA peut agir dans n'importe quel état membre, mais doit \textbf{négocier} avec les autorités locales.
	\end{itemize}
\end{shaded}

\newpage

Pour rendre ces agences \textbf{crédibles}, voici une \textbf{hiérarchie de secret} inspirée des standards militaires et du renseignement :

\begin{center}
\rotatebox{90}{
	\begin{tabular}{|>{\bfseries}P{3cm}|>{\bfseries}P{3cm}|P{5cm}|P{3cm}|}
	\hline
	\textbf{Niveau}					&	\textbf{Description}									&	\textbf{Accès}												&	\textbf{Exemple d'information classée}	\\ \hline
		Omega-1 (Top Secret +)		&	Accès ultra-restreint. Seuls 5-10 personnes par pays.	&	Chefs d'État, directeurs des services secrets.				&	Existence de l'agence, noms des Grand Anciens, localisation des artefacts maudits.	\\ \hline
		Omega-2 (Secret Défense)	&	Accès limité aux hauts gradés.							&	Generaux, directeurs de la DGSE/IntCen.						&	Missions en cours, identité des agents de terrain.	\\ \hline
		Omega-3 (Confidentiel)		&	Accès aux agents opérationnels.							&	Agents de Félicie/EORA, scientifiques habilités.			&	Protocoles d'intervention, listes des cultes surveillés.	\\ \hline
		Omega-4 (Réservé)			&	Accès aux collaborateurs externes.						&	Police, militaires en mission ponctuelle.					&	Signes de reconnaissance, protocoles de base.	\\ \hline
		Non classé					&	Information publique (désinformation).					&	Médias, grand public.										&	Rumeurs sur des "enquêtes paranormales", fausses théories du complot.	\\ \hline
	\end{tabular} }
\end{center}


\newpage

\subsection{Structure Organisationnelle}
L'EORA est organisée en divisions \textbf{multilingues} et \textbf{multinationales} :

\begin{center}
\begin{tabular}{|>{\bfseries}P{4cm}|P{7cm}|}
	\hline
	\textbf{Division}										&	\textbf{Rôle} \\ \hline
	Direction Générale Européenne (Conseil des Gardiens)	&	Supervise les opérations transnationales. Composée de hauts fonctionnaires de l'UE et d'anciens chefs des services de renseignement nationaux. \\ \hline
	Bureau des Opérations Transnationales					&	Coordonne les missions entre pays. \\ \hline
	Centre de Recherche Occulte								&	Étudie les menaces et développe des contre-mesures. Intègre des experts de diverses nationalités. \\ \hline
	Service de Contre-Espionnage et Sécurité				&	Traque les infiltrations (cultes, espions, entités). \\ \hline
	Unité de Nettoyage										&	Efface les preuves, gère les crises médiatiques à l'échelle européenne. \\ \hline
	Division Innovation et Technologie						&	Développe des outils adaptés aux menaces européennes (artefacts anciens, technologies modernes). \\ \hline
	\end{tabular}
\end{center}

\subsection{Profil des Agents}

Les agents de l'EORA doivent maîtriser \textbf{plusieurs langues} et avoir une connaissance approfondie des \textbf{cultures et mythologies locales}.

\begin{shaded}
	\textbf{Exemple de recrutement :}
	Un scientifique du CERN, ayant découvert une \textbf{anomalie dimensionnelle} lors d'une expérience, est approché par l'EORA. On lui propose de "rejoindre l'équipe pour étudier le phénomène" ou de "disparaître dans un accident de laboratoire".
\end{shaded}

\newpage

\begin{center}
\rotatebox{90}{
	\begin{tabular}{|>{\bfseries}P{3cm}|>{\bfseries}P{5cm}|P{5cm}|P{3cm}|}
	\hline
	\textbf{Type}			&	\textbf{Origine}														&	\textbf{Compétences}							&	\textbf{Faiblesses} \\ \hline
	Agent National Détaché	&	DGSE (France), \newline BND (Allemagne), \newline MI6 (Royaume-Uni)...	&	Renseignement, infiltration, combat				&	Méfiance envers les autres nationalités \\ \hline
	Scientifique Européen	&	CERN, ESA, universités européennes										&	Recherche, analyse, protocoles occultes			&	Arrogance, obsession pour la science \\ \hline
	Occultiste Repenti		&	Ancien membre de cultes européens										&	Connaissance des rituels, détection des menaces	&	Instabilité mentale, tentations \\ \hline
	Diplomate Européen		&	Commission européenne, ambassades										&	Négociation, couverture internationale			&	Manipulateur, risque de corruption \\ \hline
	Expert en Sécurité		&	Interpol, Europol, Agences Spécifiques\ldots~ 							&	Coordination internationale, gestion de crise	&	Bureaucratie, lenteur administrative \\ \hline
	\end{tabular} }
\end{center}

\newpage

\subsection{Méthodes et Outils}

\begin{itemize}
	\item \textbf{Technologies de pointe} :
	\begin{itemize}
		\item Surveillance électronique (satellites, drones).
		\item Analyse de données (IA spécialisée dans la détection de symboles occultes).
		\item Cyberdéfense (protection contre les cyberattaques occultes).
	\end{itemize}
	\item \textbf{Réseaux d'experts locaux} :
	\begin{itemize}
		\item Collaboration avec des \textbf{institutions scientifiques européennes} (CERN, ESA).
		\item Contacts dans les \textbf{universités} et les \textbf{musees} (ex : \textbf{British Museum}, \textbf{Louvre}).
	\end{itemize}
	\item \textbf{Bases de données sécurisées} :
	\begin{itemize}
		\item \textbf{Archives des phénomènes occultes européens} (vampires, loups-garous, fées).
		\item \textbf{Cartographie des lieux maudits} (forêts, châteaux, temples).
	\end{itemize}
	\item \textbf{Artefacts et Équipements} :
	\begin{itemize}
		\item Armes bénites (adaptées aux menaces locales).
		\item Détecteurs de magie (basés sur des technologies modernes et des savoirs anciens).
		\item Talismans de protection (ex : \textbf{Croix de Saint Benoît}, efficace contre les démons).
	\end{itemize}
\end{itemize}

\newpage

\subsection{Menaces Typiques}

\begin{itemize}
	\item \textbf{Organisations transnationales} :
	\begin{itemize}
		\item \textbf{L'Ordre du Dragon Noir} : Culte actif en \textbf{Allemagne, France, Italie}, cherchant à réveiller un \textbf{dragon ancien}.
		\item \textbf{La Main de Cthulhu} : Réseau lovecraftien infiltré dans les \textbf{gouvernements européens}.
		\item \textbf{Les Héritiers de Vlad} : Culte de vampires basé en \textbf{Roumanie} et \textbf{Hongrie}.
	\end{itemize}
	\item \textbf{Entités surnaturelles} :
	\begin{itemize}
		\item \textbf{Vampires} (Europe de l'Est, liés aux légendes de \textbf{Dracula}).
		\item \textbf{Loups-garous} (Allemagne, France, liées aux malédictions familiales).
		\item \textbf{Fées et Korrigans} (Irlande, Bretagne, capables de manipuler la réalité).
		\item \textbf{Démons bibliques} (liés à l'histoire chrétienne, actifs dans les \textbf{monastères abandonnés}).
	\end{itemize}
	\item \textbf{Phénomènes paranormaux} :
	\begin{itemize}
		\item \textbf{Portails dimensionnels} (ex : \textbf{Triangle des Bermudes européen} en mer du Nord).
		\item \textbf{Expériences nazies occultes} (héritage de la \textbf{Seconde Guerre mondiale}, ex : \textbf{Projet Thule}).
		\item \textbf{Lieux maudits} (ex : \textbf{Forêt de Hoia-Baciu} en Roumanie, \textbf{Château de Houska} en République tchèque).
	\end{itemize}
\end{itemize}

\newpage

\subsection{Scénarios}

\subsubsection{Opération Dragon Noir}
\textbf{Contexte :} \textbf{Allemagne, France, Italie}, années 2020.
	\begin{itemize}
		\item \textbf{Acte 1} : Des \textbf{disparitions mystérieuses} ont lieu dans plusieurs villes européennes. Les agents de l'EORA découvrent que l'\textbf{Ordre du Dragon Noir} est impliqué.
		\item \textbf{Acte 2} : L'enquête mène à un \textbf{temple caché en Forêt-Noire}, où le culte tente de réveiller un \textbf{dragon ancien}.
		\item \textbf{Acte 3} : Les agents doivent \textbf{infiltrer le rituel} et empêcher le réveil du dragon, tout en évitant les \textbf{pièges mortels} du culte.
	\item \textbf{Dilemme final} :
	\begin{itemize}
		\item Détruire le temple (et risquer de libérer le dragon).
		\item Utiliser un \textbf{artefact maudit} pour sceller le dragon (au prix de leur âme).
		\item Laisser le culte accomplir le rituel (et condamner l'Europe).
	\end{itemize}
\end{itemize}

\newpage

\subsubsection{Le Projet Thule}
\textbf{Contexte :} \textbf{Autriche, 1945-2020}.
\begin{itemize}
	\item Les agents découvrent que des \textbf{expériences nazies occultes} (Projet Thule) ont laissé derrière elles des \textbf{portails dimensionnels} instables.
	\item \textbf{Mission} : Localiser et fermer ces portails avant qu'ils ne libèrent des \textbf{entités dangereuses}.
	\item \textbf{Lieux} : \textbf{Château de Wewelsburg} (Allemagne), \textbf{bunkers alpins} (Autriche).
	\item \textbf{Menaces} : \textbf{Soldats fantômes}, \textbf{créatures lovecraftiennes}, \textbf{pièges temporels}.
\end{itemize}

\newpage

\subsubsection{La Malédiction de Hoia-Baciu}
\textbf{Contexte :} \textbf{Roumanie, Forêt de Hoia-Baciu}.
\begin{itemize}
	\item Des \textbf{phénomènes paranormaux} (disparitions, voix, lumières étranges) sont signalés dans la forêt.
	\item Les agents doivent enquêter et découvrir que la forêt est un \textbf{portail naturel} vers une autre dimension.
	\item \textbf{Dilemme} : Fermer le portail (et condamner les âmes piégées) ou le laisser ouvert (et risquer une invasion).
\end{itemize}

\newpage

%  Comparaison Félicie vs EORA vs Delta Green 
\chapter{Comparaison des Agences}

\newpage

\begin{center}
\rotatebox{90}{
	\begin{tabular}{|>{\bfseries}P{3cm}|P{4cm}|P{4cm}|P{4cm}|}
	\hline
	\textbf{Critère}				&	\textbf{Delta Green (USA)}					&	\textbf{Félicie (France)}					&	\textbf{EORA (Europe)} \\ \hline
	\textbf{Cadre Institutionnel}	&	Branche secrète du FBI/NSA					&	Ministère de la Défense / DGSE				&	IntCen (UE) / Europol \\ \hline
	\textbf{Niveau de Secret}		&	Top Secret -- Majestic-12					&	Défense Nationale -- Omega-1				&	UE Classifié -- Omega \\ \hline
	\textbf{Financement}			&	Budget noir du Pentagone					&	Fonds secrets (ministère de la Guerre)		&	Budget UE + contributions nationales \\ \hline
	\textbf{Profil des Agents}		&	Ex-militaires, ex-FBI, scientifiques		&	Légionnaires, DGSE, universitaires			&	Eurocrates, Interpol, scientifiques EU \\ \hline
	\textbf{Méthodes}				&	Force brute + technologie					&	Barbouzerie + occultisme contrôlé			&	Coopération inter-services + diplomatie occulte \\ \hline
	\textbf{Menaces Typiques}		&	Cultes lovecraftiens, extraterrestres		&	Cultes médiévaux, folklore français			&	Organisations transnationales, mythes européens \\ \hline
	\textbf{Défis Uniques}			&	Paranoïa, trahisons internes				&	Méfiance historique, cadre légal strict		&	Barrières linguistiques, conflits de juridiction \\ \hline
	\textbf{Lieux d'Opération}		&	Monde entier (focus USA)					&	France et anciennes colonies				&	Europe + pays partenaires \\ \hline
	\textbf{Alliés}					&	CIA, NSA, militaires						&	Franc-maçonnerie, Rosicruciens				&	Interpol, CERN, ESA \\ \hline
	\textbf{Faiblesses}				&	Manque de ressources, corruption interne	&	Méfiance envers l'état, manque de moyens	&	Bureaucratie, divergences culturelles \\ \hline
	\end{tabular} }
\end{center}

\newpage

%  Fiches de Personnages Prêts à Jouer 
\chapter{Fiches de Personnages}

\newpage

\section{Félicie : Jean Dupont -- Le Baroudeur}

%% \begin{shaded}
{\footnotesize 
	\begin{center}
		\textbf{Jean Dupont -- Agent de Félicie (Baroudeur)}
	\end{center}
	
	\begin{tabular}{|P{1.75cm}|P{3.75cm}|P{1.75cm}|P{3.75cm}|}
	\hline
	\textbf{Nom}			&	Jean Dupont						&	\textbf{Âge}			&	38 ans \\ \hline
	\textbf{Origine}		&	Légion étrangère (13e DBLE) 	&	\textbf{Grade}			&	Capitaine (officiellement "en congé sans solde") \\ \hline
	\textbf{Apparence}		&	\multicolumn{3}{ p{10cm}|}{1m85, cheveux courts, cicatrice sur la joue gauche, yeux gris. Porte toujours une veste en cuir usée.} \\ \hline
	\textbf{Compétences}	&	
		Combat rapproché (Couteau, Pistolet) : 85\%~\newline
		Tir (Fusil d'assaut, Pistolet silencieux) : 80\%~\newline
		Infiltration : 75\%~\newline
		Survie (Désert, Jungle) : 90\%~\newline
		Résistance à la folie : 70\%~\newline
		Langues : Français (100\%), Anglais (60\%), Arabe (40\%)~\newline
	& \textbf{Équipement}	&	
		Pistolet silencieux (Heckler \& Koch MK 23) + 12 balles bénites~\newline
		Couteau de combat (lame en fer froid)~\newline
		Gilet pare-balles (marqué "Police Nationale")~\newline
		Détecteur de magie (artefact ancien)~\newline
		Carte des bases secrètes de Félicie~\newline
		Faux passeport (nom : "Pierre Laurent")
	\\ \hline
	\textbf{Personnalité}	&	
		Calme sous la pression, mais sujet à des crises de paranoïa. -- 
		Loyal envers Félicie, mais commence à douter de ses méthodes. -- 
		Aime le whisky et les cigares cubains. -- 
		Hanté par les missions passées (a perdu des amis en opération). 
	& \textbf{Faiblesses}	&	
		Paranoïa (doit faire un jet de SAN toutes les 24h en mission).~\newline
		Alcoolisme (boit pour oublier).~\newline
		Phobie des araignées (à cause d'une mission au Cambodge).
	\\ \hline
	\textbf{Historique}	&	
		Ancien légionnaire, a servi en Afrique et au Moyen-Orient.~\newline
		Recruté par Félicie après avoir vu un "fantôme" en mission au Yémen.~\newline
		A participé à l'opération "Ethor bedi zure erresuma" (a perdu 30 ans de mémoire). 
	& & \\ \hline
	\end{tabular} }
%% \end{shaded}

\newpage

\section{Félicie : Marie Laurent -- Le Scientifique}

%% \begin{shaded}
{\footnotesize 
	\begin{center}
	\textbf{Marie Laurent -- Agent de Félicie (Scientifique)}
	\end{center}

	\begin{tabular}{|P{1.75cm}|P{3.75cm}|P{1.75cm}|P{3.75cm}|}
	\hline
	\textbf{Nom}			&	Marie Laurent 									&	\textbf{Âge}			&	42 ans \\ \hline
	\textbf{Origine}		&	CNRS (physicienne spécialisée en énergie noire)	&	\textbf{Grade}			&	Directrice de recherche (officiellement "en congé sabbatique") \\ \hline
	\textbf{Apparence}		&	\multicolumn{3}{ p{10cm}|}{1m70, cheveux bruns en chignon, lunettes à monture d'or, toujours vêtue de blanc. } \\ \hline
	\textbf{Compétences}	&	
		Physique quantique : 95\%~\newline
		Occultisme : 80\%~\newline
		Chimie (analyse d'artefacts) : 90\%~\newline
		Médecine (premiers secours) : 70\%~\newline
		Langues : Français (100\%), Anglais (90\%), Latin (60\%), Allemand (50\%)
	& \textbf{Équipement}	&	
		Ordinateur portable (avec base de données sur les artefacts)~\newline
		Spectromètre portable (détecte les énergies surnaturelles)~\newline
		Grimoire protégé ("De Vermis Mysteriis", édition française)~\newline
		Trousse de premiers secours + sédatifs~\newline
		Pistolet (pour la défense personnelle)
	\\ \hline
	\textbf{Personnalité}	&	
		Obsédée par la compréhension des phénomènes occultes.~\newline
		Méthodique et logique, mais peut devenir fanatique.~\newline
		A un sens de l'humour noir.~\newline
		Se méfie des militaires (les trouve "trop brutaux").
	& \textbf{Faiblesses}	&	
		Obsession pour le savoir interdit (doit faire un jet de SAN si elle lit un grimoire maudit).~\newline
		Phobie des insectes (à cause d'une expérience avec un artefact maudit).~\newline
		Insomnie chronique.
	\\ \hline
	\textbf{Historique}	&	
		A découvert un "portail dimensionnel" dans un laboratoire du CNRS.~\newline
		Recrutée par Félicie après avoir tenté de publier ses découvertes.~\newline
		A travaillé sur le projet "Janus" (confinement d'un Grand Ancien).
	& & \\ \hline
	\end{tabular} }
%% \end{shaded}

\newpage

\section{EORA : Klaus Schmidt -- Le Scientifique Européen}

%% \begin{shaded}
{\footnotesize 
	\begin{center}
		\textbf{Klaus Schmidt -- Agent de l'EORA (Scientifique)}
	\end{center}
	
	\begin{tabular}{|P{1.75cm}|P{3.75cm}|P{1.75cm}|P{3.75cm}|}
	\hline
	\textbf{Nom}			&	Klaus Schmidt 					&	\textbf{Âge}			&	45 ans \\ \hline
	\textbf{Origine}		&	CERN (physicien des particules)	&	\textbf{Grade}			&	Chef de projet (officiellement "en mission temporaire") \\ \hline
	\textbf{Apparence}		&	\multicolumn{3}{ p{10cm}|}{1m78, cheveux gris, barbe taillée, porte toujours une blouse blanche. } \\ \hline
	\textbf{Compétences}	&	
		Physique des particules : 95\%~\newline
		Informatique (hacking, analyse de données) : 85\%~\newline
		Occultisme (théorique) : 70\%~\newline
		Langues : Allemand (100\%), Anglais (100\%), Français (80\%), Italien (50\%)
	& \textbf{Équipement}	&	
		Ordinateur portable (accès aux bases de données de l'EORA)~\newline
		Détecteur de radiations (modifié pour détecter les énergies occultes)~\newline
		Badge d'accès aux installations du CERN~\newline
		Pistolet (pour la défense personnelle)
	\\ \hline
	\textbf{Personnalité}	&	
		Rationnel et sceptique (mais a vu trop de choses pour nier l'occulte).~\newline
		Perfectionniste, parfois arrogant.~\newline
		Aime la musique classique (Bach, Beethoven).~\newline
		Se méfie des politiques.
	& \textbf{Faiblesses}	&	
		Arrogance (méprise les "non-scientifiques").
		Obsession pour les "preuves tangibles" (peut mettre en danger la mission).~\newline
		Phobie des espaces confinés (à cause d'un incident au CERN).~\newline
	\\ \hline
	\textbf{Historique}	&	
		A découvert une "anomalie dimensionnelle" lors d'une expérience au CERN.~\newline
		Recruté par l'EORA après avoir tenté de publier ses résultats.~\newline
		A travaillé sur le projet "Thule" (fermeture des portails nazis).
	& & \\ \hline
	\end{tabular} }
%% \end{shaded}

\newpage

\section{EORA : Sofia Petrov -- L'Agent de Terrain}

%% \begin{shaded}
{\footnotesize 
	\begin{center}
		\textbf{Sofia Petrov -- Agent de l'EORA (Terrain)}
	\end{center}
	
	\begin{tabular}{|P{1.75cm}|P{3.75cm}|P{1.75cm}|P{3.75cm}|}
	\hline	
	\textbf{Nom}			&	Sofia Petrov								&	\textbf{Âge}			&	32 ans \\ \hline
	\textbf{Origine}		&	BND (Allemagne) / Ex-Spetsnaz (Russie) 		&	\textbf{Grade}			&	Agent spécial (officiellement "consultante en sécurité") \\ \hline
	\textbf{Apparence}		&	\multicolumn{3}{ p{10cm}|}{1m75, cheveux blonds courts, yeux bleus, tatouage discret (une croix orthodoxe). } \\ \hline
	\textbf{Compétences}	&	
		Combat rapproché : 90\%~\newline
		Tir (tous types) : 85\%~\newline
		Infiltration : 80\%~\newline
		Langues : Russe (100\%), Allemand (100\%), Anglais (90\%), Français (70\%), Polonais (50\%)
	& \textbf{Équipement}	&	
		Pistolet (Glock 17) + silencieux~\newline
		Couteau de combat (lame en argent)~\newline
		Tenue de combat noire (marquée "Interpol")~\newline
		Détecteur de mensonges (basé sur des technologies avancées)~\newline
		Passeport diplomatique (faux)
	\\ \hline
	\textbf{Personnalité}	&	
		Froid et calculatrice, mais loyale envers ses coéquipiers.~\newline
		A un sens aigu de la justice.~\newline
		Déteste les nazis (son grand-père a combattu sur le front de l'Est).~\newline
		Superstitieuse (porte toujours une médaille de Saint Georges).
	& \textbf{Faiblesses}	&	
		Méfiance envers les occultistes (les considère comme "dangereux").~\newline
		Problèmes de confiance (a été trahie par un coéquipier).~\newline
		Phobie des serpents.
	\\ \hline
	\textbf{Historique}	&	
		Ancien agent du BND, spécialisée dans les opérations en Europe de l'Est.~\newline
		Recrutée par l'EORA après avoir enquêté sur un culte en Roumanie.~\newline
		A participé à l'opération "Dragon Noir".
	& & \\ \hline
	\end{tabular} }
%% \end{shaded}

\newpage

%  Idées de Scénarios 
\chapter{Idées de Scénarios}

\section{Scénarios pour Félicie}

\subsection{Le Musée Maudit (Paris, 2020)}
\begin{shaded}
\textbf{Niveau :} Débutant à Intermédiaire \\ \textbf{Durée :} 2-3 séances \\ \textbf{Lieux :} Paris, Catacombes, Château de Chantilly

\textbf{Synopsis :}
Des objets maudits sont volés dans plusieurs musées parisiens. Les agents de Félicie sont chargés de les retrouver avant qu'ils ne soient utilisés pour un rituel visant à réveiller une entité ancienne.

\textbf{Déroulement :}
	\begin{enumerate}
	\item \textbf{Enquête initiale} : Les agents sont envoyés au \textbf{Musée du Louvre}, où un \textbf{artefact égyptien} (un scarabée maudit) a été volé. Ils découvrent que d'autres objets ont disparu dans d'autres musées (\textbf{Musée de Cluny}, \textbf{Musée d'Orsay}).
	\item \textbf{Piste du voleur} : Les vols semblent liés à un \textbf{ancien conservateur} du Louvre, \textbf{Henri de Montclair}, disparu il y a 10 ans. Il aurait rejoint un culte occulte.
	\item \textbf{Infiltration du culte} : Les agents traquent le culte jusqu'aux \textbf{catacombes de Paris}, où ils découvrent un \textbf{autel dédié à Chaugnar Faugn}. Le culte prépare un rituel pour \textbf{libérer une entité}.
	\item \textbf{Affrontement final} : Dans les catacombes, les agents doivent affronter le culte et l'entité invoquée (un \textbf{Shoggoth} ou un \textbf{Mi-Go}).
	\item \textbf{Dilemme} :
	\begin{itemize}
		\item Détruire l'autel (et risquer de libérer l'entité).
		\item Utiliser un artefact pour "sceller" l'entité (mais au prix de leur santé mentale).
	\item Laisser le culte accomplir le rituel (et condamner Paris).
	\end{itemize}
\end{enumerate}

\textbf{PNJ Importants :}
	\begin{itemize}
		\item \textbf{Henri de Montclair} : Ancien conservateur, maintenant \textbf{Grand Prêtre} du culte. Charismatique et manipulateur.
		\item \textbf{élodie Moreau} : Journaliste enquêtant sur les vols. Peut aider les agents (ou les trahir).
		\item \textbf{Inspecteur Dubois} : Policier chargé de l'enquête officielle. Sceptique, mais peut être convaincu.
	\end{itemize}

\textbf{Menaces :}
	\begin{itemize}
		\item \textbf{Cultistes} (10-15, armés de couteaux et de pistolets).
		\item \textbf{Garde du culte} (2-3, anciens mercenaires, armés de fusils d'assaut).
		\item \textbf{Shoggoth} (1, dans les catacombes).
	\end{itemize}

\textbf{Récompenses :}
	\begin{itemize}
		\item Récupération des artefacts volés.
		\item Information sur un \textbf{autre culte} (pour un scénario futur).
		\item Possible promotion pour les agents.
	\end{itemize}
\end{shaded}

\newpage

\subsection{La Forêt des Ombres (Bretagne, 2020)}
\begin{shaded}
\textbf{Niveau :} Intermédiaire \\ \textbf{Durée :} 3-4 séances \\ \textbf{Lieux :} Forêt de Brocéliande, Village de Tréhorenteuc

\textbf{Synopsis :}
Des \textbf{disparitions mystérieuses} ont lieu dans la forêt de Brocéliande. Les agents de Félicie découvrent que les \textbf{korrigans} et les \textbf{fées} locales sont impliqués, mais quelque chose de bien plus dangereux rôde dans les bois...

\textbf{Déroulement :}
\begin{enumerate}
	\item \textbf{Enquête initiale} : Les agents sont envoyés dans le village de \textbf{Tréhorenteuc}, où plusieurs habitants ont disparu. Les villageois parlent de "lumières dans la forêt" et de "voix qui appellent".
	\item \textbf{Rencontre avec les korrigans} : En explorant la forêt, les agents rencontrent des \textbf{korrigans}, qui leur offrent de l'aide... en échange d'un service. (\textbf{Attention} : les korrigans sont des manipulateurs !)
	\item \textbf{Découverte du portails} : Les agents trouvent un \textbf{cercle de pierres} (un portails vers les \textbf{Contrées du Rêve}). Le culte local (les \textbf{Enfants de la Forêt}) tente de l'utiliser pour invoquer une entité.
	\item \textbf{Affrontement final} : Les agents doivent affronter le culte et une \textbf{créature des Contrées du Rêve} (un \textbf{Nightgaunt} ou un \textbf{Hound of Tindalos}).
	\item \textbf{Dilemme} :
	\begin{itemize}
		\item Détruire le portail (et condamner les âmes piégées dans la forêt).
		\item Utiliser le portail pour voyager dans les Contrées du Rêve (risque : ne jamais revenir).
		\item Laisser le culte accomplir le rituel (et risquer une invasion).
	\end{itemize}
\end{enumerate}

\textbf{PNJ Importants :}
	\begin{itemize}
		\item \textbf{Morgane Le Goff} : Vieille femme du village, connaît les légendes locales. Peut guider les agents (ou les induire en erreur).
		\item \textbf{Yann le Rouge} : Chef du culte des \textbf{Enfants de la Forêt}. Fanatique et dangereux.
		\item \textbf{Le Korrigan} : Être malicieux qui propose son aide... pour un prix.
	\end{itemize}

\textbf{Menaces :}
	\begin{itemize}
		\item \textbf{Korrigans} (5-6, utilisent des illusions).
		\item \textbf{Cultistes} (8-10, armés de couteaux et de haches).
		\item \textbf{Nightgaunt} (1, dans les Contrées du Rêve).
	\end{itemize}

\textbf{Récompenses :}
	\begin{itemize}
		\item Découverte d'un \textbf{artefact féerique} (peut être utile pour un scénario futur).
		\item Information sur les \textbf{Contrées du Rêve}.
	\end{itemize}
\end{shaded}

\newpage

\section{Scénarios pour l'EORA}

\subsection{Opération Dragon Noir (Europe, 2020)}
\begin{shaded}
\textbf{Niveau :} Avancé \\ \textbf{Durée :} 4-5 séances \\ \textbf{Lieux :} Berlin, Forêt-Noire, Prague

\textbf{Synopsis :}
L'\textbf{Ordre du Dragon Noir}, un culte transnational, tente de réveiller un dragon ancien sous la Forêt-Noire. Les agents de l'EORA doivent les arrêter avant qu'il ne soit trop tard.

\textbf{Déroulement :}
\begin{enumerate}
	\item \textbf{Enquête initiale} : Des \textbf{disparitions} et des \textbf{symboles étranges} (runes) sont signalés à Berlin, Prague, et dans la Forêt-Noire. Les agents de l'EORA sont envoyés pour enquêter.
	\item \textbf{Infiltration du culte} : Les agents découvrent que l'Ordre du Dragon Noir est dirigé par un \textbf{ancien officier nazi}, \textbf{Reinhard von Drachenfels}, qui a survécu à la guerre grâce à un rituel de longévité.
	\item \textbf{Découverte du temple} : Dans la Forêt-Noire, les agents trouvent un \textbf{temple caché} où le culte prépare le rituel de réveil.
	\item \textbf{Affrontement final} : Les agents doivent affronter les \textbf{gardes du culte} et \textbf{von Drachenfels} lui-même, qui tente d'accomplir le rituel.
	\item \textbf{Dilemme final} :
	\begin{itemize}
		\item Détruire le temple (et risquer de libérer le dragon).
		\item Utiliser un \textbf{artefact maudit} (le "C\oe ur de Dragon") pour sceller le dragon (au prix de leur âme).
		\item Laisser le culte accomplir le rituel (et condamner l'Europe).
	\end{itemize}
\end{enumerate}

\textbf{PNJ Importants :}
	\begin{itemize}
		\item \textbf{Reinhard von Drachenfels} : Ancien officier nazi, chef du culte. Charismatique et impitoyable.
		\item \textbf{Anna Müller} : Scientifique allemande, ancienne membre du culte. Peut aider les agents (ou les trahir).
		\item \textbf{Inspecteur Kovacs} : Policier hongrois, sceptique mais compétent.
	\end{itemize}

\textbf{Menaces :}
	\begin{itemize}
		\item \textbf{Gardes du culte} (10-12, armés de fusils d'assaut et de lances-flammes).
		\item \textbf{Reinhard von Drachenfels} (1, utilise des sorts noirs).
		\item \textbf{Dragon ancien} (1, si le rituel est accompli).
	\end{itemize}

\textbf{Récompenses :}
	\begin{itemize}
		\item Récupération du "Cœur de Dragon" (artefact puissant).
		\item Information sur d'autres cultes transnationaux.
	\end{itemize}
\end{shaded}

\newpage

\subsection{Le Projet Thule (Autriche, 2020)}
\begin{shaded}
\textbf{Niveau :} Avancé \\ \textbf{Durée :} 3-4 séances \\ \textbf{Lieux :} Château de Wewelsburg (Allemagne), Bunkers alpins (Autriche)

\textbf{Synopsis :}
Des \textbf{phénomènes paranormaux} (disparitions, voix, apparitions de soldats fantômes) sont signalés dans les \textbf{Alpes autrichiennes}. Les agents de l'EORA découvrent que des \textbf{expériences nazies occultes} (Projet Thule) ont laissé derrière elles des \textbf{portails dimensionnels} instables.

\textbf{Déroulement :}
	\begin{enumerate}
	\item \textbf{Enquête initiale} : Les agents sont envoyés dans un village autrichien où des habitants ont disparu. Ils découvrent des \textbf{symboles nazis} et des \textbf{runes anciennes} gravées dans les montagnes.
	\item \textbf{Découverte du bunker} : Les agents trouvent l'entrée d'un \textbf{bunker nazi abandonné}, où des expériences occultes ont été menées.
	\item \textbf{Exploration du bunker} : À l'intérieur, ils découvrent des \textbf{portails dimensionnels} et des \textbf{soldats fantômes} (les esprits des scientifiques nazis).
	\item \textbf{Affrontement final} : Les agents doivent fermer les portails avant qu'ils ne libèrent des \textbf{entités dangereuses} (ex : un \textbf{Shoggoth} ou un \textbf{Byakhee}).
	\item \textbf{Dilemme} :
	\begin{itemize}
		\item Fermer les portails (et condamner les âmes piégées à l'intérieur).
		\item Étudier les portails (risque : invasion).
		\item Détruire le bunker (et perdre une source de connaissances occultes).
	\end{itemize}
	\end{enumerate}

\textbf{PNJ Importants :}
	\begin{itemize}
		\item \textbf{Docteur Ernst Vogel} : Ancien scientifique nazi, maintenant un \textbf{fantôme}. Peut aider les agents (ou les induire en erreur).
		\item \textbf{Capitaine Hans Bauer} : Ancien officier nazi, toujours vivant grâce à un rituel. Fanatique et dangereux.
		\item \textbf{Inspecteur Weber} : Policier autrichien, sceptique mais déterminé.
	\end{itemize}

\textbf{Menaces :}
	\begin{itemize}
		\item \textbf{Soldats fantômes} (5-6, armés de fusils de la Seconde Guerre mondiale).
		\item \textbf{Capitaine Hans Bauer} (1, utilise des sorts noirs).
		\item \textbf{Shoggoth} (1, si un portail est ouvert).
	\end{itemize}

\textbf{Récompenses :}
	\begin{itemize}
		\item Récupération de \textbf{documents nazis occultes} (peut être utile pour un scénario futur).
		\item Information sur le \textbf{Projet Thule} et ses ramifications.
	\end{itemize}
\end{shaded}

\newpage

%  Localisations et Descriptions 
\chapter{Localisations et Descriptions}

\section{France}

\subsection{Paris}
\begin{itemize}
	\item \textbf{Lieux clés} :
	\begin{itemize}
		\item \textbf{Catacombes de Paris} : Réseau de tunnels cachant des secrets occultes. \textbf{Ambiance} : Sombre, humide, rempli d'ossements. \textbf{Menaces} : Cultes, fantômes, créatures des ombres.
		\item \textbf{Musée du Louvre} : Abrite des artefacts maudits (ex : \textbf{Le Miroir de la Duchesse de Montpensier}). \textbf{Ambiance} : Luxueux, mais avec des zones interdites (sous-sols, réserves).
		\item \textbf{Quartier Latin} : Siège de sociétés secrètes (ex : \textbf{Loge maçonnique occulte}). \textbf{Ambiance} : étudiants, librairies anciennes, ruelles sombres.
		\item \textbf{Bibliothèque nationale de France} : Contient des grimoires protégés. \textbf{Ambiance} : Silencieuse, labyrinthique, gardée par des bibliothécaires méfiants.
	\end{itemize}
\end{itemize}

\subsection{Forêt de Brocéliande (Bretagne)}
\begin{itemize}
	\item \textbf{Description} : Forêt légendaire liée aux légendes arthuriennes. \textbf{Ambiance} : Mystérieuse, brumeuse, avec des arbres centenaires.
	\item \textbf{Lieux clés} :
	\begin{itemize}
		\item \textbf{Val sans Retour} : Lieu où les \textbf{fées} attirent les voyageurs.
		\item \textbf{Tombeau de Merlin} : Site occulte où des rituels ont lieu.
		\item \textbf{Château de Comper} : Ruines d'un château lié à la légende de Lancelot.
		\end{itemize}
	\item \textbf{Menaces} : Korrigans, fées, loups-garous, cultes païens.
\end{itemize}

\subsection{Château de Brissac (Maine-et-Loire)}
\begin{itemize}
	\item \textbf{Description} : Château hanté, surnommé le "Géant Vert" à cause de ses tours imposantes. \textbf{Ambiance} : Gothique, sombre, avec des bruits étranges la nuit.
	\item \textbf{Légende} : Le château serait hanté par le fantôme de \textbf{Jacques de Brézé}, assassiné par son épouse et son amant.
	\item \textbf{Menaces} : Fantômes, poltergeists, cultes sataniques.
\end{itemize}

\newpage

\section{Europe}

\subsection{Forêt-Noire (Allemagne)}
\begin{itemize}
	\item \textbf{Description} : Forêt dense et mystérieuse, liée aux légendes germaniques. \textbf{Ambiance} : Sombre, avec des arbres géants et des brumes épaisses.
	\item \textbf{Lieux clés} :
	\begin{itemize}
		\item \textbf{Temple du Dragon} : Lieu de rituels anciens, caché sous une colline.
		\item \textbf{Lac Noir} : Lac maudit où des disparitions ont lieu.
		\item \textbf{Village de Schwarzwald} : Village isolé où les habitants pratiquent des cultes païens.
	\end{itemize}
	\item \textbf{Menaces} : Loups-garous, dragons, cultes occultes, esprits de la forêt.
\end{itemize}

\subsection{Château de Wewelsburg (Allemagne)}
\begin{itemize}
	\item \textbf{Description} : Château utilisé par les nazis pour des rituels occultes. \textbf{Ambiance} : Sinistre, avec des symboles nazis et des runes anciennes.
	\item \textbf{Légende} : Le château aurait été le centre du \textbf{Projet Thule}, visant à utiliser des forces occultes pour gagner la guerre.
	\item \textbf{Menaces} : Soldats fantômes, entités lovecraftiennes, pièges magiques.
\end{itemize}

\subsection{Forêt de Hoia-Baciu (Roumanie)}
\begin{itemize}
	\item \textbf{Description} : Forêt connue comme la "Bermude Triangle de la Transylvanie". \textbf{Ambiance} : Effrayante, avec des arbres tordus et des zones où le temps semble s'arrêter.
	\item \textbf{Légende} : Des \textbf{portails dimensionnels} et des \textbf{disparitions} y sont signalés depuis des siècles.
	\item \textbf{Menaces} : Vampires, loups-garous, entités des Contrées du Rêve, cultes locaux.
\end{itemize}

\newpage

\subsection{Château de Houska (République tchèque)}
\begin{itemize}
	\item \textbf{Description} : Château médiéval construit sur un "portail vers l'enfer". \textbf{Ambiance} : Gothique, avec des murs épais et des sous-sols interdits.
	\item \textbf{Légende} : Le château aurait été construit pour "sceller" une faille dimensionnelle.
	\item \textbf{Menaces} : Démons, esprits, cultes sataniques.
\end{itemize}

\newpage

%  Idées Diverses 
\chapter{Idées Diverses}

\section{Adaptation des Règles de Delta Green}

\begin{itemize}
	\item \textbf{Utilisation des compétences} :
	\begin{itemize}
		\item Les agents de Félicie et de l'EORA peuvent utiliser les \textbf{règles de Delta Green} pour les jets de compétences.
		\item \textbf{Compétences spécifiques} :
		\begin{itemize}
			\item \textbf{Occultisme} : Pour reconnaître les symboles, rituels, ou entités.
			\item \textbf{Histoire de l'occulte} : Pour comprendre le contexte historique des menaces.
			\item \textbf{Langues} : Indispensable pour les agents de l'EORA (plusieurs langues européennes).
		\end{itemize}
	\end{itemize}
	\item \textbf{Équipement} :
	\begin{itemize}
		\item Les agents peuvent utiliser des \textbf{armes bénites} (dégâts supplémentaires contre les créatures surnaturelles).
		\item Les \textbf{artefacts} peuvent avoir des effets spéciaux (ex : un grimoire peut donner un bonus en Occultisme, mais au prix de la santé mentale).
	\end{itemize}
	\item \textbf{Santé mentale} :
	\begin{itemize}
		\item Les agents de Félicie et de l'EORA doivent faire des \textbf{jets de SAN} (Santé Mentale) lorsqu'ils sont confrontés à des phénomènes occultes.
		\item La \textbf{perte de SAN} peut entraîner des \textbf{troubles psychologiques} (paranoïa, hallucinations, etc.).
	\end{itemize}
	\item \textbf{Rituels et Magie} :
	\begin{itemize}
		\item Les agents peuvent apprendre des \textbf{rituels de protection} ou de \textbf{neutralisation} (ex : exorcisme, cercle de protection).
		\item L'utilisation de la magie noire est \textbf{interdite} (sauf en dernier recours), et peut entraîner des \textbf{conséquences graves} (corruption, possession).
	\end{itemize}
\end{itemize}

\section{Création de Personnages}

\begin{itemize}
	\item \textbf{Origines possibles} :
	\begin{itemize}
		\item \textbf{Félicie} : Légionnaire, policier, scientifique, occultiste repenti, diplomate.
		\item \textbf{EORA} : Agent de renseignement (DGSE, BND, MI6), scientifique (CERN, ESA), expert en sécurité (Interpol, Europol).
	\end{itemize}
	\item \textbf{Compétences recommandées} :
	\begin{itemize}
		\item \textbf{Combat} : Pour les agents de terrain.
		\item \textbf{Occultisme} : Pour comprendre les menaces.
		\item \textbf{Langues} : Indispensable pour l'EORA.
		\item \textbf{Psychologie} : Pour résister à la folie.
	\end{itemize}
	\item \textbf{Équipement de base} :
	\begin{itemize}
		\item Arme de service (pistolet ou fusil).
		\item Équipement de protection (gilet pare-balles, talisman).
		\item Outils d'enquête (ordinateur portable, détecteur de magie).
	\end{itemize}
	\item \textbf{Traits de personnalité} :
	\begin{itemize}
		\item Chaque personnage doit avoir au moins un \textbf{défaut} (paranoïa, obsession, phobie) et une \textbf{qualité} (loyauté, intelligence, courage).
	\end{itemize}
\end{itemize}

\newpage

\section{Conseils pour le Maître de Jeu}

\begin{itemize}
	\item \textbf{Ambiance} :
	\begin{itemize}
		\item Utilisez des \textbf{descriptions détaillées} pour immerger les joueurs (bruits, odeurs, atmosphère).
		\item Jouez sur la \textbf{peur psychologique} (doute, trahison, folie).
	\end{itemize}
	\item \textbf{Gestion des PNB} :
	\begin{itemize}
		\item Les \textbf{cultistes} et les \textbf{entités} doivent être \textbf{menaçants} et \textbf{imprévisibles}.
		\item Utilisez des \textbf{PNJ ambivalents} (allies potentiels, mais avec des secrets).
	\end{itemize}
	\item \textbf{Scénarios} :
	\begin{itemize}
		\item Prévoyez des \textbf{dilemmes moraux} (sauver un innocent ou accomplir la mission ?).
		\item Utilisez des \textbf{twists} (trahisons, révélations choquantes).
	\end{itemize}
	\item \textbf{Règles maison} :
	\begin{itemize}
		\item Vous pouvez ajouter des \textbf{règles spécifiques} pour Félicie ou l'EORA (ex : bonus pour les agents français en France, malus pour les barrières linguistiques en Europe).
	\end{itemize}
\end{itemize}

\newpage

\section{Inspirations pour de Futurs Scénarios}

\begin{itemize}
	\item \textbf{Cultes historiques} :
	\begin{itemize}
		\item \textbf{Les Cathares} : Culte médiéval persécuté, mais toujours actif dans le sud de la France.
		\item \textbf{Les Templiers} : Ordre militaire lié à des secrets occultes (ex : \textbf{le Saint-Graal}).
		\item \textbf{La Franc-Maçonnerie occulte} : Sociétés secrètes avec des rituels interdits.
	\end{itemize}
	\item \textbf{Légendes locales} :
	\begin{itemize}
		\item \textbf{La Bête du Gévaudan} : Créature mystérieuse qui pourrait être un \textbf{loup-garou} ou un \textbf{démon}.
		\item \textbf{Les Marie-Morgane} : Sirènes maléfiques de Bretagne.
		\item \textbf{Le Juif Errant} : Légende liée à des malédictions éternelles.
	\end{itemize}
	\item \textbf{événements historiques} :
	\begin{itemize}
		\item \textbf{La Révolution française} : Période de chaos où des cultes occultes ont prospéré.
		\item \textbf{La Seconde Guerre mondiale} : Expériences nazies occultes (ex : \textbf{Projet Thule}).
		\item \textbf{La Guerre froide} : Opérations secrètes liées à l'occulte (ex : \textbf{Projet MKUltra}).
	\end{itemize}
	\item \textbf{Lieux mystérieux} :
	\begin{itemize}
		\item \textbf{Le Mont-Saint-Michel} : Abbaye liée à des légendes de \textbf{dragons} et de \textbf{fantômes}.
		\item \textbf{Les Alignements de Carnac} : Site mégalithique avec des propriétés occultes.
		\item \textbf{Le Lac Pavin} : Lac maudit en Auvergne, lié à des \textbf{sacrifices} anciens.
	\end{itemize}
\end{itemize}

\newpage

% =============================================
% CHAPITRE : SCÉNARIOS POUR MUSÉES EUROPÉENS
% =============================================

\chapter{Scénarios pour Musées Européens}
\markboth{Scénarios pour Musées Européens}{}

\section*{Introduction aux Scénarios}
\addcontentsline{toc}{section}{Introduction aux Scénarios}

Les musées européens regorgent d'\textbf{artefacts maudits}, de \textbf{cultes secrets} et de \textbf{phénomènes inexpliqués}, en faisant des lieux idéaux pour des aventures horrifiques dans l'univers de \textbf{Delta Green}, \textbf{Félicie} ou \textbf{EORA}.
Ce chapitre propose des scénarios détaillés pour des musées emblématiques de \textbf{Paris, Berlin, Londres, Madrid, Amsterdam} et \textbf{Saint-Pétersbourg}, avec des menaces adaptées à chaque lieu.

\begin{shaded}
\textbf{Note :} Ces scénarios peuvent être joués indépendamment ou liés pour former une campagne plus large (ex : une conspiration internationale impliquant plusieurs musées).
\end{shaded}


\section{Scénarios pour Paris}

\newpage

\subsection{Musée National d'Histoire Naturelle (MNHN) -- \textit{L'Éveil des Oubliés}}
\markright{L'Éveil des Oubliés}

\begin{shaded}
	\textbf{Niveau :} Intermédiaire \quad
	\textbf{Durée :} 3-4 séances \quad
	\textbf{Lieux :} Galerie de Paléontologie, Cabinet des Curiosités, Laboratoires souterrains
\end{shaded}

\subsubsection*{Synopsis}
Une exposition temporaire sur les \textit{Espèces disparues et légendes} au MNHN est le théâtre de \textbf{disparitions mystérieuses} parmi le personnel et les visiteurs. Les agents découvrent que les ossements exposés (notamment ceux d'un \textit{dragon médiéval} et d'une \textit{créature des cavernes}) ne sont pas de simples fossiles\ldots~  mais des \textbf{restes de Grands Anciens} ou d'entités scellées par des rituels anciens. Un \textbf{scientifique renégat} (ancien membre de Félicie) tente de \textbf{réveiller ces entités}.

\subsubsection*{Déroulement}
\begin{center}
	\begin{tabular}{|P{3cm}|P{9cm}|}
	\hline
	\textbf{Acte} 							& \textbf{Événements clés} \\ \hline
	\textbf{1. L'Appel du MNHN} 			& Les agents sont contactés par un conservateur paniqué (\textbf{Dr. Élise Moreau}). Des gardiens ont disparu près de la Galerie de Paléontologie, et des visiteurs rapportent avoir vu des \textit{ombres bouger} la nuit. \\ \hline
	\textbf{2. Le Cabinet des Horreurs} 	& Les agents explorent le \textbf{Cabinet des Curiosités} (fermé au public depuis 1940). Ils y découvrent un \textbf{grimoire} (\textit{Codex Bestiarum}) décrivant un rituel : \textit{réveiller les os}. \\ \hline
	\textbf{3. Les Laboratoires Interdits} 	& Dans les sous-sols, les agents trouvent des \textbf{cages vides} (marquées \textit{Spécimen X-47}) et des notes sur des expériences menées dans les années 1930 par un savant nazi (lié à l'\textbf{Ahnenerbe}). \\ \hline
	\textbf{4. Le Rituel Final} 			& Le \textbf{Dr. René Lefèvre} (ancien de Félicie, maintenant corrompu) tente de compléter le rituel dans la \textbf{Grande Galerie de l'Évolution}. \\ \hline
	\end{tabular}
\end{center}

\subsubsection*{Dilemme final}
\begin{itemize}
	\item \textbf{Détruire le grimoire} (et perdre une connaissance précieuse).
	\item \textbf{Utiliser le grimoire} pour contrôler les entités (risque : corruption).
\end{itemize}

\subsubsection*{Récompenses}
\begin{itemize}
	\item \textbf{Artefact :} Le \textit{Codex Bestiarum} (bonus en Occultisme, corruption utilisateur).
	\item \textbf{Allié :} Sophie Lambert peut devenir un contact.
	\item \textbf{Information :} Lien entre le MNHN et l'Ahnenerbe.
\end{itemize}

\subsubsection*{Conseils pour le MJ}
\begin{itemize}
	\item \textbf{Ambiance :} Bruits de pas dans les galeries vides, chuchotements en latin.
	\item \textbf{Twist :} Le \textit{dragon} du musée est un Grand Ancien endormi (un \textbf{Yig} ou un \textbf{Dhole}).
	\item \textbf{Variante :} Si les joueurs échouent, le MNHN devient un repaire de monstres (\textit{scénario futur : Le Musée Maudit}).
\end{itemize}

\newpage

\subsubsection*{PNJ Importants}
\begin{center}
	\begin{tabular}{|P{2cm}|P{3cm}|P{3.5cm}|P{2.5cm}|} 
	\hline
	\textbf{Nom} 		& \textbf{Rôle} 			& \textbf{Description} & \textbf{Motivation} \\ \hline
	Dr. Élise Moreau 	& Conservatrice 			& 50 ans, spécialiste des cétacés préhistoriques. A remarqué les disparitions croire. 	& Peur : Veut protéger le musée. \\ \hline
	Dr. René Lefèvre 	& Ancien agent de Félicie 	& 60 ans, pâle, yeux jaunes. Porte une bague en os. 							& Folie : Croit que les entités peuvent \textit{sauver la science}. \\ \hline
	Lucien Dubois 		& Gardien de nuit 			& 30 ans, superstitieux. A vu des \textit{choses} mais ne veut pas en parler. 			& Lâcheté : Prêt à fuir. \\ \hline
	Sophie Lambert 		& Étudiante en archéologie 	& 25 ans, curieuse. A trouvé un fragment du grimoire. 						& Ambition : Veut publier sa découverte. \\ \hline
	\end{tabular}
\end{center}

\subsubsection*{Menaces}
\begin{center}
	\begin{tabular}{|P{2cm}|P{3cm}|P{3.5cm}|P{2.5cm}|}
	\hline
	\textbf{Type} 			& \textbf{Description} 							& \textbf{Statistiques (DG)} 									& \textbf{Faiblesses} \\ \hline
	Os Vivants 				& Squelettes animés par une énergie maléfique. 	& FOR 70\%, DEX 60\%, CON 80\% (immunes aux balles normales). 	& Armes bénites, feu. \\ \hline
	Ombre de X-47 			& Entité amorphe (ancien spécimen du musée). 	& FOR 90\%, DEX 80\%, CON 100\% (invisible dans le noir). 		& Lumière vive, sels bénits. \\ \hline
	Dr. Lefèvre (corrompu) 	& Utilise des sorts noirs. 						& FOR 50\%, DEX 70\%, CON 60\%, Magie 80\%. 					& Détruire le grimoire. \\ \hline
	\end{tabular}
\end{center}

\newpage

\subsection{Musée du Louvre -- \textit{Le Souffle de la Joconde}}

\markright{Le Souffle de la Joconde}

\begin{shaded}
	\textbf{Niveau :} Débutant à Intermédiaire \quad
	\textbf{Durée :} 2-3 séances \quad
	\textbf{Lieux :} Galerie Denon, Sous-sols interdits, Atelier de restauration
\end{shaded}

\subsubsection*{Synopsis}
La \textbf{Joconde} a disparu pendant une nuit. Officiellement, c'est un vol, mais les agents de Félicie découvrent que le tableau a été \textbf{remplacé par une copie}\ldots~  et que ceux qui l'ont touchée tombent dans un \textbf{coma inexplicable}. Pire : les caméras montrent des \textbf{ombres} se déplaçant autour du tableau, comme si quelque chose en sortait.

\subsubsection*{Déroulement}
\begin{center}
	\begin{tabular}{|P{3cm}|P{9cm}|}
	\hline
	\textbf{Acte} 						& \textbf{Événements clés} \\ \hline
	\textbf{1. L'Enquête Officielle}	& Les agents sont envoyés sous couverture de \textit{consultants en sécurité}. La police pense à un vol organisé. \\ \hline
	\textbf{2. La Copie Maudite}		& En analysant la copie, les agents découvrent qu'elle a été peinte au 16ème siècle avec des \textbf{pigments inconnus} (sang + poudre d'os). \\ \hline
	\textbf{3. Les Sous-Sols Oubliés}	& Dans les sous-sols du Louvre se trouve un \textbf{atelier secret} où des restaurateurs ont recréé des œuvres maudites. \\ \hline
	\textbf{4. La Femme Sans Yeux}		& Dans la Salle des Caryatides, les agents affrontent l'entité : une version démoniaque de Mona Lisa. \\ \hline
	\end{tabular}
\end{center}

\subsubsection*{Dilemme final}

\begin{itemize}
	\item \textbf{Détruire la copie} (et perdre un chef-d'œuvre).
	\item \textbf{Sceller le démon} (mais un agent devra rester en sacrifice).
\end{itemize}

\subsubsection*{Récompenses}
\begin{itemize}
	\item \textbf{Artefact :} La vraie Joconde (peut servir de talisman contre les illusions).
	\item \textbf{Allié :} Claire Delacroix peut devenir un contact.
	\item \textbf{Information :} Réseau de collectionneurs occultes (Comte de Montclair).
\end{itemize}

\newpage

\subsubsection*{PNJ Importants}
\begin{center}
	\begin{tabular}{|P{2cm}|P{3cm}|P{3.5cm}|P{2.5cm}|}
	\hline
	\textbf{Nom} 		& \textbf{Rôle} 			& \textbf{Description}									& \textbf{Motivation} \\ \hline
	Commissaire Bernard & Policier 					& 45 ans, sceptique. Croit à un vol classique. 			& Devoir : Résolution rapide de l'affaire. \\ \hline
	Claire Delacroix 	& Restauratrice 			& 35 ans, perfectionniste. A travaillé sur la copie. 	& Culpabilité : Sait que quelque chose ne va pas. \\ \hline
	Marcel Moreau 		& Gardien (en coma) 		& 60 ans. Rêve de la Femme Sans Yeux. 					& Désespoir : Hurle et saigne des yeux s'il se réveille. \\ \hline
	Comte de Montclair 	& Collectionneur occulte 	& 70 ans, riche, mystérieux. A commandé la copie. 		& Obsession : Posséder la vraie Joconde. \\ \hline
	\end{tabular}
\end{center}

\subsubsection*{Menaces}

\begin{center}
	\begin{tabular}{|P{2cm}|P{3cm}|P{3.5cm}|P{2.5cm}|}
	\hline
	\textbf{Type} 		& \textbf{Description} 									& \textbf{Statistiques (DG)} 							& \textbf{Faiblesses} \\ \hline
	Statues Animées 	& Vénus de Milo, Victoire de Samothrace\ldots~  prennent vie. 	& FOR 80\%, DEX 50\%, CON 90\% (immunes aux balles). 	& Armes en fer, eau bénite. \\ \hline
	Femme Sans Yeux 	& Entité liée au tableau. Peut hypnotiser (jet de VOL ou perte de 1D6 SAN). 	& FOR 60\%, DEX 90\%, CON 100\%, Magie 85\% 	& Miroir, nom de Rê. \\ \hline
	Ombres du Louvre 	& Esprits des anciens voleurs/visiteurs piégés. 		& FOR 40\%, DEX 80\%, CON 50\% 							& Lumière vive, exorcisme. \\ \hline
	\end{tabular}
\end{center}

\newpage

\subsection{Musée de l'Homme -- \textit{Les Masques de la Folie}}
\markright{Les Masques de la Folie}

\begin{shaded}
	\textbf{Niveau :} Avancé \quad
	\textbf{Durée :} 4-5 séances \quad
	\textbf{Lieux :} Galerie des Crânes, Réserves anthropologiques, Laboratoire de génétique
\end{shaded}

\subsubsection*{Synopsis}

Le Musée de l'Homme organise une exposition sur les \textit{Masques Rituels du Monde}. Depuis l'arrivée d'un \textbf{masque dogon}, des visiteurs et employés souffrent de \textbf{hallucinations collectives} : ils voient des \textbf{visages déformés} et entendent des \textbf{voix dans leur tête}. Certains disparaissent, et leurs corps réapparaissent plus tard\ldots~  \textbf{sans visage}.

\subsubsection*{Déroulement}
\begin{center}
	\begin{tabular}{|P{3cm}|P{9cm}|}
	\hline
	\textbf{Acte} 							& \textbf{Événements clés} \\ \hline
	\textbf{1. Les Premiers Signes} 		& Les agents sont appelés après qu'un gardien a \textit{enlevé son propre visage} avec ses mains. \\ \hline
	\textbf{2. Le Masque Maudite} 			& Le masque est scellé dans une vitrine blindée, mais quiconque le regarde doit faire un jet de SAN (perte de 1D10). \\ \hline
	\textbf{3. Les Réserves Interdites} 	& Les agents explorent les réserves, où des masques maudits sont entreposés. Ils y trouvent des notes du Dr. Diallo. \\ \hline
	\textbf{4. Le Rituel du Visage Vide} 	& Dans la Galerie des Crânes, les agents doivent affronter le Démon des Masques. \\ \hline
	\end{tabular}
\end{center}

\subsubsection*{Dilemme final}
\begin{itemize}
	\item \textbf{Détruire le masque} (et libérer les âmes piégées).
	\item \textbf{Sceller le masque à nouveau} (mais un agent devra porter le masque).
\end{itemize}

\newpage

\subsubsection*{PNJ Importants}
\begin{center}
	\begin{tabular}{|P{2cm}|P{3cm}|P{3.5cm}|P{2.5cm}|}
	\hline
	\textbf{Nom} 			& \textbf{Rôle} 	& \textbf{Description} 															& \textbf{Motivation} \\ \hline
	Dr. Amadou Diallo 		& Anthropologue		& 40 ans, spécialiste des masques africains. A disparu en étudiant le masque. 	& Curiosité : Comprendre le rituel. \\ \hline
	Cécile Moreau 			& Conservatrice 	& 30 ans, stressée. A vu des choses mais refuse d'en parler. 					& Peur : Peur de perdre son visage. \\ \hline
	Le Gardien Sans Visage 	& Victime 			& Ancien gardien. Erre dans le musée, cherchant son visage. 					& Désespoir : Attaque quiconque s'approche. \\ \hline
	\end{tabular}
\end{center}

\subsubsection*{Menaces}
\begin{center}
	\begin{tabular}{|P{2cm}|P{3cm}|P{3.5cm}|P{2.5cm}|}
	\hline
	\textbf{Type} 		& \textbf{Description} 									& \textbf{Statistiques (DG)} 					& \textbf{Faiblesses} \\ \hline
	Masques Animés 		& Masque Baoulé, Masque Fang\ldots~  possèdent les visiteurs. 	& FOR 70\%, DEX 60\%, CON 80\% 					& Détruire le masque physique, exorcisme. \\ \hline
	Démon des Masques 	& Entité amorphe qui vole les visages. 					& FOR 90\%, DEX 80\%, CON 100\%, Magie 90\% 	& Miroir brisé, sang de son gardien. \\ \hline
	Visiteurs Possédés 	& Employés ou touristes contrôlés par les masques. 		& FOR 60\%, DEX 70\%, CON 50\% 					& Dégâts normaux (libère l'âme). \\ \hline
	\end{tabular}
\end{center}

\newpage

\subsection{Musée du Quai Branly -- \textit{Les Tambours de la Mort}}
\markright{Les Tambours de la Mort}


\begin{shaded}
	\textbf{Niveau :} Avancé \quad
	\textbf{Durée :} 3-4 séances \quad
	\textbf{Lieux :} Galerie des Arts Africains, Réserves des Instruments, Jardin des Colonnes
\end{shaded}

\subsubsection*{Synopsis}

Une exposition sur les \textit{Instruments Sacrés d'Afrique} est le théâtre de morts mystérieuses : des visiteurs meurent \textbf{en dansant}, comme possédés, leurs yeux explosés et leurs oreilles saignantes. Les agents découvrent que les tambours exposés sont maudits : ils \textbf{appellent les esprits des ancêtres}\ldots~  qui veulent se venger.

\subsubsection*{Déroulement}
\begin{center}
	\begin{tabular}{|P{3cm}|P{9cm}|}
	\hline
	\textbf{Acte} 						& \textbf{Événements clés} \\ \hline
	\textbf{1. La Danse Macabre} 		& Les agents sont appelés après qu'un groupe de touristes a été retrouvé mort, figés en position de danse. \\ \hline
	\textbf{2. Les Tambours Maudits} 	& Les agents découvrent que les tambours émettent des fréquences inaudibles qui contrôlent l'esprit. \\ \hline
	\textbf{3. Le Rituel Interrompu} 	& Dans les réserves, les agents trouvent un autel avec des offrandes fraîches (sang, plumes). \\ \hline
	\textbf{4. Le Dernier Tambour} 		& Les agents doivent détruire le tambour principal avant qu'il ne réveille un Loa. \\ \hline
	\end{tabular}
\end{center}

\subsubsection*{Dilemme final}

\begin{itemize}
	\item \textbf{Détruire le tambour} (et perdre un artefact historique).
	\item \textbf{Le garder} (mais risquer une malédiction permanente).
\end{itemize}

\newpage

\subsubsection*{PNJ Importants}
\begin{center}
	\begin{tabular}{|P{2cm}|P{3cm}|P{3.5cm}|P{2.5cm}|}
	\hline
	\textbf{Nom} 		& \textbf{Rôle} 	& \textbf{Description} 										& \textbf{Motivation} \\ \hline
	Papa Legba 			& Prêtre vaudou 	& 60 ans, charismatique, porte un chapeau en plumes. 		& Vengeance : Veut libérer les esprits. \\ \hline
	Dr. Anna Kowalski 	& Ethnologue 		& 35 ans, passionnée. A étudié les tambours. 				& Culpabilité : Sait que les tambours sont dangereux. \\ \hline
	Thomas Leroy 		& Gardien 			& 40 ans, superstitieux. A entendu les tambours la nuit. 	& Peur : Refuse de monter la garde. \\ \hline
	\end{tabular}
\end{center}

\subsubsection*{Menaces}
\begin{center}
	\begin{tabular}{|P{2cm}|P{3cm}|P{3.5cm}|P{2.5cm}|}
	\hline
	\textbf{Type} 		& \textbf{Description} 									& \textbf{Statistiques (DG)} 	& \textbf{Faiblesses} \\ \hline
	Esprits Ancestraux 	& Invisibles, possèdent les vivants. 					& FOR 80\%, DEX 70\%, CON 60\% 	& Gri-gri, exorcisme vaudou. \\ \hline
	Tambours Animés 	& Battent tout seuls, émettent des ondes maléfiques. 	& -- 							& Détruire le tambour principal. \\ \hline
	Possédés 			& Visiteurs ou employés contrôlés par les esprits. 		& FOR 70\%, DEX 60\%, CON 50\% 	& Dégâts normaux (libère l'âme). \\ \hline
	\end{tabular}
\end{center}

\newpage

\section{Scénarios pour Berlin}
\subsection{Pergamonmuseum -- \textit{La Porte des Dieux Oubliés}}
\markright{La Porte des Dieux Oubliés}

\begin{shaded}
	\textbf{Niveau :} Avancé \quad
	\textbf{Durée :} 4-5 séances \quad
	\textbf{Lieux :} Porte d'Ishtar, Salle du Trône de Babylone, Sous-sols archéologiques
\end{shaded}

\subsubsection*{Synopsis}
La \textbf{Porte d'Ishtar} (reconstruite au Pergamonmuseum) s'est mise à saigner une \textbf{substance noire} pendant la nuit. Les archéologues qui l'ont touchée sont tombés dans un coma, et leurs rêves décrivent une \textbf{cité perdue} où règne un \textbf{dieu-araignée}. Les agents de l'EORA découvrent que la porte est en réalité un \textbf{portail scellé}\ldots~  et que quelqu'un (un culte néonazi) tente de l'ouvrir.

\subsubsection*{Déroulement}
\begin{center}
	\begin{tabular}{|P{3cm}|P{9cm}|}
	\hline
	\textbf{Acte} 						& \textbf{Événements clés} \\ \hline
	\textbf{1. Le Sang sur la Porte} 	& Les agents sont envoyés après qu'un gardien a disparu près de la Porte d'Ishtar. Des traces de substance noire sont retrouvées. \\ \hline
	\textbf{2. Le Culte de l'Araignée} 	& Les agents découvrent qu'un groupe néonazi (les \textit{Enfants de Tiamat}) utilise le musée pour réveiller un dieu ancien. \\ \hline
	\textbf{3. Les Sous-Sols Maudits} 	& Sous le musée, les agents trouvent une salle secrète où le culte a reconstruit un autel babylonien. \\ \hline
	\textbf{4. L'ouverture du Portail} 	& Les agents doivent empêcher Hans Weber de compléter le rituel avant que la Porte d'Ishtar ne s'ouvre. \\ \hline
	\end{tabular}
\end{center}

\subsubsection*{Dilemme final}
\begin{itemize}
\item \textbf{Détruire la porte} (et perdre un chef-d'œuvre historique).
\item \textbf{Sceller le portail} (mais un agent devra rester en sacrifice).
\end{itemize}

\newpage

\subsubsection*{PNJ Importants}
\begin{center}
	\begin{tabular}{|P{2cm}|P{3cm}|P{3.5cm}|P{2.5cm}|}
	\hline
	\textbf{Nom} 		& \textbf{Rôle} 	& \textbf{Description} 											& \textbf{Motivation} \\ \hline
	Hans Weber 			& Chef du culte 	& 50 ans, ancien archéologue. Porte une bague en obsidienne. 	& Folie : Croit que les dieux méritent de régner. \\ \hline
	Dr. Clara Schmidt 	& Conservatrice 	& 40 ans, experte en Mésopotamie. A remarqué les anomalies. 	& Peur : Veut protéger le musée. \\ \hline
	Jürgen Müller		& Gardien 			& 30 ans, nerveux. A vu des choses dans les sous-sols. 			& Désespoir : Prêt à fuir. \\ \hline
	\end{tabular}
\end{center}

\subsubsection*{Menaces}
\begin{center}
	\begin{tabular}{|P{2cm}|P{3cm}|P{3.5cm}|P{2.5cm}|}
	\hline
	\textbf{Type} 		& \textbf{Description} 											& \textbf{Statistiques (DG)} 		& \textbf{Faiblesses} \\ \hline
	Statues Animées 	& Tiamat, Kingu\ldots~  attaquent avec des pouvoirs divins. 			& FOR 90\%, DEX 50\%, CON 100\% 	& Armes bénites, eau du Tigre/Euphrate. \\ \hline
	Cultistes Néonazis 	& Armés (fusils, couteaux). Utilisent des symboles occultes. 	& FOR 60\%, DEX 70\%, CON 50\% 		& Dégâts normaux. \\ \hline
	Araignée de Tiamat 	& Géante, invisible (sauf pour les fous). 						& FOR 100\%, DEX 80\%, CON 90\% 	& Feu, symboles de Marduk. \\ \hline
	\end{tabular}
\end{center}

\newpage

\subsection{Neues Museum -- \textit{Le Buste qui Chuchote}}
\markright{Le Buste qui Chuchote}

\begin{shaded}
	\textbf{Niveau :} Intermédiaire \quad
	\textbf{Durée :} 2-3 séances \quad
	\textbf{Lieux :} Salle de Néfertiti, Laboratoire de restauration, Sous-sols égyptiens
\end{shaded}

\subsubsection*{Synopsis}

Le \textbf{buste de Néfertiti} a commencé à parler aux visiteurs. Ceux qui l'ont entendu tombent dans un état catatonique, répétant sans cesse : *"Elle attend sous le sable"*. Les agents découvrent que le buste est lié à un rituel égyptien pour réveiller une entité\ldots~  et qu'un archéologue corrompu tente de l'utiliser.

\subsubsection*{Déroulement}
\begin{center}
	\begin{tabular}{|P{3cm}|P{9cm}|}
	\hline
	\textbf{Acte} 						& \textbf{Événements clés} \\ \hline
	\textbf{1. Les Voix du Musée} 		& Les agents sont appelés après qu'un groupe de touristes a été retrouvé catatonique devant le buste. \\ \hline
	\textbf{2. Le Secret de Néfertiti} 	& Les agents découvrent que le buste est maudit : il appelle les esprits des pharaons morts. \\ \hline
	\textbf{3. Le Laboratoire Maudite} 	& Dans le laboratoire de restauration, les agents trouvent des notes du Dr. Bauer : il a volé un artefact (un Scarabée du Nil). \\ \hline
	\textbf{4. Le Réveil de la Reine} 	& Les agents doivent empêcher le Dr. Bauer de réveiller Néfertiti (qui n'est pas une reine\ldots~  mais un démon). \\ \hline
	\end{tabular}
\end{center}

\subsubsection*{Dilemme final}
\begin{itemize}
	\item \textbf{Détruire le buste} (et perdre un artefact historique).
	\item \textbf{Sceller le démon} (mais un agent devra porter le scarabée maudit).
\end{itemize}

\newpage

\subsubsection*{PNJ Importants}
\begin{center}
	\begin{tabular}{|P{2cm}|P{3cm}|P{3.5cm}|P{2.5cm}|}
	\hline
	\textbf{Nom} 		& \textbf{Rôle} 		& \textbf{Description} 									& \textbf{Motivation} \\ \hline
	Dr. Karl Bauer 		& Archéologue corrompu 	& 55 ans, obsédé. Porte un collier en or. 				& Ambition : Veut devenir immortel. \\ \hline
	Dr. Amina Hassan 	& Égyptologue 			& 30 ans, intelligente. A remarqué les anomalies. 		& Devoir : Veut protéger l'héritage égyptien. \\ \hline
	Markus Vogel 		& Gardien 				& 40 ans, superstitieux. A entendu les chuchotements. 	& Peur : Refuse de s'approcher du buste. \\ \hline
	\end{tabular}
\end{center}

\subsubsection*{Menaces}
\begin{center}
	\begin{tabular}{|P{2cm}|P{3cm}|P{3.5cm}|P{2.5cm}|}
	\hline
	\textbf{Type} 			& \textbf{Description} 							& \textbf{Statistiques (DG)} 					& \textbf{Faiblesses} \\ \hline
	Momies Vivantes 		& Squelettiques, rapides. 						& FOR 70\%, DEX 80\%, CON 60\% 					& Feu, eau bénite. \\ \hline
	Dr. Bauer (corrompu) 	& Utilise des sorts égyptiens. 					& FOR 50\%, DEX 70\%, CON 60\%, Magie 80\% 		& Détruire le scarabée. \\ \hline
	Néfertiti (Démon) 		& Hypnotise (jet de VOL ou perte de 1D10 SAN). 	& FOR 80\%, DEX 90\%, CON 100\%, Magie 95\% 	& Miroir, nom de Rê. \\ \hline
	\end{tabular}
\end{center}

\newpage

\subsection{DDR Museum -- \textit{Les Archives de l'Oubli}}
\markright{Les Archives de l'Oubli}

\begin{shaded}
	\textbf{Niveau :} Débutant à Intermédiaire \quad
	\textbf{Durée :} 2-3 séances \quad
	\textbf{Lieux :} Salle des Espions, Archives secrètes, Sous-sol du musée
\end{shaded}

\subsubsection*{Synopsis}
Le DDR Museum organise une exposition sur les \textit{Secrets de la Stasi}. Mais des visiteurs rapportent avoir vu des \textbf{hommes en noir} qui n'existent pas sur les photos. Pire : certains disparaissent, et leurs dossiers (dans les archives du musée) s'effacent comme par magie. Les agents découvrent que la Stasi a mené des expériences occultes\ldots~  et que quelque chose est revenu.

\subsubsection*{Déroulement}
\begin{center}
	\begin{tabular}{|P{3cm}|P{9cm}|}
	\hline
	\textbf{Acte} 	& \textbf{Événements clés} \\ \hline
	\textbf{1. Les Hommes en Noir} 		& Les agents sont envoyés après qu'un journaliste a disparu en enquêtant sur la Stasi. \\ \hline
	\textbf{2. Les Archives Maudites} 	& Les agents découvrent que la Stasi a testé des rituels pour créer des espions invisibles. \\ \hline
	\textbf{3. Le Sous-Sol Interdit} 	& Sous le musée, les agents trouvent une salle d'interrogatoire où des prisonniers ont été sacrifiés. \\ \hline
	\textbf{4. Le Projet Schattenspiel} & Les agents doivent empêcher Klaus Schmidt de compléter le rituel. \\ \hline
	\end{tabular}
\end{center}

\subsubsection*{Dilemme final}
\begin{itemize}
	\item \textbf{Détruire les archives} (et perdre des preuves historiques).
	\item \textbf{Sceller le portail} (mais un agent devra porter un uniforme de la Stasi maudit).
\end{itemize}

\newpage

\subsubsection*{PNJ Importants}
\begin{center}
	\begin{tabular}{|P{2cm}|P{3cm}|P{3.5cm}|P{2.5cm}|}
	\hline
	\textbf{Nom} 		& \textbf{Rôle} 			& \textbf{Description} 										& \textbf{Motivation} \\ \hline
	Klaus Schmidt 		& Ancien agent de la Stasi 	& 60 ans, froid, calculateur. Porte un uniforme vintage. 	& Nostalgie : Veut restaurer la RDA. \\ \hline
	Anna Weber 			& Journaliste 				& 30 ans, courageuse. A découvert le projet. 				& Vérité : Veut révéler les crimes de la Stasi. \\ \hline
	Inspecteur Müller 	& Police 					& 45 ans, sceptique. Pense à un canular. 					& Devoir : Veut résoudre l'affaire. \\ \hline
	\end{tabular}
\end{center}

\subsubsection*{Menaces}
\begin{center}
	\begin{tabular}{|P{2cm}|P{3cm}|P{3.5cm}|P{2.5cm}|}
	\hline
	\textbf{Type} 				& \textbf{Description} 								& \textbf{Statistiques (DG)} 				& \textbf{Faiblesses} \\ \hline
	Espions Fantômes 			& Invisibles (sauf avec une caméra à infrarouge). 	& FOR 60\%, DEX 90\%, CON 50\% 				& Lumière vive, symboles de la RFA. \\ \hline
	Klaus Schmidt (corrompu) 	& Utilise des pouvoirs de la Stasi. 				& FOR 50\%, DEX 70\%, CON 60\%, Magie 75\% 	& Détruire son uniforme. \\ \hline
Ombre Vivante 				& Entité qui absorbe la lumière. 					& FOR 80\%, DEX 80\%, CON 90\% 				& Lumière blanche, miroirs. \\ \hline
	\end{tabular}
\end{center}

\newpage
	
\section{Scénarios pour Autres Capitales Européennes}
\subsection{British Museum (Londres) -- \textit{Le Sceau de Salomon}}
\markright{Le Sceau de Salomon}

\begin{shaded}
	\textbf{Niveau :} Avancé \quad
	\textbf{Durée :} 4 séances \quad
	\textbf{Lieux :} Salle des Trésors, Réserves juives, Sous-sols maudits
\end{shaded}

\subsubsection*{Synopsis}
Un \textbf{sceau en or} (supposé être celui du Roi Salomon) a été volé dans la salle des Trésors Juifs. Depuis, des employés du musée disparaissent, et leurs corps sont retrouvés \textbf{mumifiés en quelques heures}. Les agents découvrent que le sceau ne protège pas\ldots~  mais \textbf{emprisonne un démon}\ldots~  et qu'un rabbin corrompu tente de le libérer.

\subsubsection*{Déroulement}
\begin{center}
	\begin{tabular}{|P{3cm}|P{9cm}|}
	\hline
	\textbf{Acte} 							& \textbf{Événements clés} \\ \hline
	\textbf{1. La Malédiction du Sceau} 	& Les agents sont appelés après la disparition d'un conservateur. Les caméras montrent une ombre ailée près du sceau. \\ \hline
	\textbf{2. Le Rabbin Hérétique}			& Les agents découvrent que le Rabin Eliézer a volé le sceau pour libérer un démon (un Asmodée). \\ \hline
	\textbf{3. Les Catacombes de Londres} 	& Le rabbin a caché le sceau dans les catacombes sous le musée. \\ \hline
	\textbf{4. Le Rituel de Libération} 	& Les agents doivent empêcher le rabbin de compléter le rituel avant que le démon ne soit libéré. \\ \hline
	\end{tabular}
\end{center}

\subsubsection*{Dilemme final}
\begin{itemize}
	\item \textbf{Détruire le sceau} (et libérer le démon).
	\item \textbf{Le sceller à nouveau} (mais un agent devra porter le sceau maudit).
\end{itemize}

\newpage

\subsubsection*{PNJ Importants}
\begin{center}
	\begin{tabular}{|P{2cm}|P{3cm}|P{3.5cm}|P{2.5cm}|}
	\hline
	\textbf{Nom} 			& \textbf{Rôle} 	& \textbf{Description} 										& \textbf{Motivation} \\ \hline
	Rabin Eliézer 			& Rabbin corrompu 	& 70 ans, vieux, sage en apparence. Porte une kippah noire. & Folie : Croit que le démon peut sauver Israël. \\ \hline
	Dr. Rachel Cohen 		& Historienne 		& 40 ans, experte en judaïsme. A remarqué les anomalies. 	& Devoir : Veut protéger le patrimoine juif. \\ \hline
	Inspecteur Thompson 	& Police 			& 50 ans, sceptique. Pense à un meurtre rituel. 			& Devoir : Veut résoudre l'affaire. \\ \hline
	\end{tabular}
\end{center}

\subsubsection*{Menaces}
\begin{center}
	\begin{tabular}{|P{2cm}|P{3cm}|P{3.5cm}|P{2.5cm}|}
	\hline
	\textbf{Type} 				& \textbf{Description} 								& \textbf{Statistiques (DG)} 				& \textbf{Faiblesses} \\ \hline
	Golems de Chair 			& Créés à partir des victimes.						& FOR 90\%, DEX 40\%, CON 80\% 				& Feu, sels bénits. \\ \hline
	Asmodée (Démon) 			& Démon de la colère. Peut lancer des malédictions. & FOR 100\%, DEX 80\%, CON 90\%, Magie 95\% & Sceau de Salomon, nom de Dieu (en hébreu). \\ \hline
	Rabin Eliézer (corrompu) 	& Utilise des pouvoirs cabalistiques. 				& FOR 40\%, DEX 60\%, CON 50\%, Magie 90\% 	& Détruire le sceau. \\ \hline
	\end{tabular}
\end{center}

\newpage

\subsection{Museo del Prado (Madrid) --~\newline \textit{Le Tableau qui Regarde}}
\markright{Le Tableau qui Regarde}

\begin{shaded}
	\textbf{Niveau :} Intermédiaire \quad
	\textbf{Durée :} 3 séances \quad
	\textbf{Lieux :} Salle de Goya, Réserves, Atelier de restauration
\end{shaded}

\subsubsection*{Synopsis}
Un tableau de Goya (*"Le Sommeil de la Raison"*) a commencé à changer : les visages des personnages bougent, et ceux qui le regardent trop longtemps tombent dans un coma. Les agents découvrent que le tableau est maudit : il \textbf{emprisonne les âmes} de ceux qui le regardent\ldots~  et qu'un peintre fou tente de libérer ces âmes.

\subsubsection*{Déroulement}
\begin{center}
	\begin{tabular}{|P{3cm}|P{9cm}|}
	\hline
	\textbf{Acte} 							& \textbf{Événements clés} \\ \hline
	\textbf{1. Les Victimes du Prado} 		& Les agents sont appelés après que des visiteurs sont tombés dans le coma devant le tableau. \\ \hline
	\textbf{2. Le Secret de Goya} 			& Les agents découvrent que Goya a utilisé des pigments maudits (mélange de sang et de poudre d'os). \\ \hline
	\textbf{3. L'Atelier Maudite} 			& Dans les réserves, les agents trouvent un atelier secret où un peintre fou recrée le tableau. \\ \hline
	\textbf{4. Le Rituel de Libération} 	& Les agents doivent empêcher Diego López de compléter le rituel. \\ \hline
	\end{tabular}
\end{center}

\subsubsection*{Dilemme final}
\begin{itemize}
	\item \textbf{Détruire le tableau} (et perdre un chef-d'œuvre).
	\item \textbf{Sceller les âmes} (mais un agent devra entrer dans le tableau).
\end{itemize}

\newpage

\subsubsection*{PNJ Importants}
\begin{center}
	\begin{tabular}{|P{2cm}|P{3cm}|P{3.5cm}|P{2.5cm}|}
	\hline
	\textbf{Nom} 			& \textbf{Rôle} 		& \textbf{Description} 									& \textbf{Motivation} \\ \hline
	Diego López 			& Peintre fou 			& 40 ans, désordonné, yeux injectés de sang. Porte une palette de couleurs noires. 	& Folie : Veut libérer les âmes\ldots~  pour les peindre. \\ \hline
	Dr. Isabella Garcia 	& Historienne de l'art 	& 35 ans, passionnée. A remarqué les changements. 		& Devoir : Veut protéger l'art espagnol. \\ \hline
	Inspecteur Ruiz 		& Police 				& 50 ans, sceptique. Pense à une intoxication. 			& Devoir : Veut résoudre l'affaire. \\ \hline
	\end{tabular}
\end{center}

\subsubsection*{Menaces}
\begin{center}
	\begin{tabular}{|P{2cm}|P{3cm}|P{3.5cm}|P{2.5cm}|}
	\hline
	\textbf{Type} 			& \textbf{Description} 									& \textbf{Statistiques (DG)} 				& \textbf{Faiblesses} \\ \hline
	Créatures Peintes 		& Monstres de Goya (ex : Saturno devorando a su hijo). 	& FOR 70\%, DEX 80\%, CON 60\% 				& Détruire le tableau, feu. \\ \hline
	Diego López (corrompu) 	& Utilise des pouvoirs artistiques. 					& FOR 40\%, DEX 70\%, CON 50\%, Magie 80\% 	& Détruire sa palette. \\ \hline
	Âmes Piégées 			& Esprits des victimes. Possèdent les vivants. 			& FOR 50\%, DEX 70\%, CON 40\% 				& Exorcisme, détruire le tableau. \\ \hline
	\end{tabular}
\end{center}

\newpage

\subsection{Rijksmuseum (Amsterdam) -- \newline \textit{Le Portrait du Diable}}
\markright{Le Portrait du Diable}


\begin{shaded}
	\textbf{Niveau :} Débutant \quad
	\textbf{Durée :} 2 séances \quad
	\textbf{Lieux :} Galerie d'Honneur, Réserves, Atelier de Rembrandt
\end{shaded}

\subsubsection*{Synopsis}
Un portrait de Rembrandt (*"L'Homme au Turban d'Or"*) a commencé à parler aux visiteurs. Ceux qui l'ont écouté sont retrouvés morts, avec un sourire figé et des yeux noirs. Les agents découvrent que le portrait est maudit : il \textbf{appelle le Diable}\ldots~  et qu'un collectionneur tente de l'utiliser.

\subsubsection*{Déroulement}
\begin{center}
	\begin{tabular}{|P{3cm}|P{9cm}|}
	\hline
	\textbf{Acte} 							& \textbf{Événements clés} \\ \hline
	\textbf{1. Les Morts Souriantes} 		& Les agents sont appelés après que des visiteurs sont morts devant le portrait. \\ \hline
	\textbf{2. Le Collectionneur Maudite} 	& Les agents découvrent que le Comte van der Meer a volé le portrait pour invoquer le Diable. \\ \hline
	\textbf{3. L'Atelier de Rembrandt} 		& Dans les réserves, les agents trouvent un atelier secret où le comte a reproduit le portrait. \\ \hline
	\textbf{4. Le Contrat Diabolique} 		& Les agents doivent empêcher le comte de signer le contrat avant que le Diable n'apparaisse. \\ \hline
	\end{tabular}
\end{center}

\subsubsection*{Dilemme final}

\begin{itemize}
	\item \textbf{Détruire le portrait} (et perdre un chef-d'œuvre).
	\item \textbf{Sceller le Diable} (mais un agent devra signer le contrat à sa place).
\end{itemize}

\newpage

\subsubsection*{PNJ Importants}
\begin{center}
	\begin{tabular}{|P{2cm}|P{3cm}|P{3.5cm}|P{2.5cm}|}
	\hline
	\textbf{Nom} 		& \textbf{Rôle} 	& \textbf{Description} 									& \textbf{Motivation} \\ \hline
	Comte van der Meer 	& Collectionneur 	& 60 ans, riche, arrogant. Porte une bague en or noir. 	& Ambition : Veut devenir immortel. \\ \hline
	Dr. Eva de Vries 	& Historienne 		& 30 ans, intelligente. A remarqué les changements. 	& Devoir : Veut protéger l'art néerlandais. \\ \hline
	Inspecteur Janssen 	& Police 			& 45 ans, sceptique. Pense à un empoisonnement. 		& Devoir : Veut résoudre l'affaire. \\ \hline
	\end{tabular}
\end{center}

\subsubsection*{Menaces}

\begin{center}
	\begin{tabular}{|P{2cm}|P{3cm}|P{3.5cm}|P{2.5cm}|}
	\hline
	\textbf{Type} 					& \textbf{Description} 				& \textbf{Statistiques (DG)} 					& \textbf{Faiblesses} \\ \hline
	Démons Mineurs 					& Serviteurs du Diable. 			& FOR 70\%, DEX 80\%, CON 60\% 					& Eau bénite, symboles religieux. \\ \hline
	Comte van der Meer (corrompu) 	& Utilise des pouvoirs diaboliques. & FOR 50\%, DEX 70\%, CON 60\%, Magie 85\% 		& Détruire le contrat. \\ \hline
	Le Diable 						& Prend la forme de Rembrandt. 		& FOR 100\%, DEX 90\%, CON 100\%, Magie 100\% 	& Foi absolue, détruire le portrait. \\ \hline
	\end{tabular}
\end{center}

\newpage

\subsection{Musée de l'Ermitage (Saint-Pétersbourg) -- \newline \textit{La Malédiction des Tsars}}
\markright{La Malédiction des Tsars}

\begin{shaded}
	\textbf{Niveau :} Avancé \quad
	\textbf{Durée :} 4 séances \quad
	\textbf{Lieux :} Salle des Trônes, Réserves impériales, Sous-sols maudits
\end{shaded}

\subsubsection*{Synopsis}
Un \textbf{\oe uf de Fabergé} (celui de 1916, disparu depuis la Révolution) a réapparu dans les collections de l'Ermitage. Depuis, des employés rapportent avoir vu des \textbf{fantômes en uniforme impérial}\ldots~  et certains disparaissent, leurs corps retrouvés \textbf{gelés} (alors que la température est normale). Les agents découvrent que l'\oe uf ne contient pas un joyau\ldots~  mais \textbf{l'âme d'un tsar maudit}\ldots~  et qu'un noble russe tente de le libérer.

\subsubsection*{Déroulement}
\begin{center}
	\begin{tabular}{|P{3cm}|P{9cm}|}
	\hline
	\textbf{Acte} 								& \textbf{Événements clés} \\ \hline
	\textbf{1. Les Fantômes de l'Ermitage} 		& Les agents sont appelés après que des gardiens ont disparu. Les caméras montrent des ombres en uniforme du 19ème siècle. \\ \hline
	\textbf{2. L'\oe uf Maudit} 				& Les agents découvrent que l'\oe uf contient une âme : celle du Tsar Nicolas II (ou d'un tsar maudit). \\ \hline
	\textbf{3. Les Archives Impériales} 		& Dans les réserves, les agents trouvent des documents sur un rituel pour libérer l'âme du tsar. \\ \hline
	\textbf{4. Le Rituel de la Restauration} 	& Les agents doivent empêcher le Comte Orlov de compléter le rituel. \\ \hline
	\end{tabular}
\end{center}

\subsubsection*{Dilemme final}
\begin{itemize}
	\item \textbf{Détruire l'\oe uf} (et perdre un trésor historique).
	\item \textbf{Sceller l'âme} (mais un agent devra porter l'\oe uf maudit).
\end{itemize}

\newpage

\subsubsection*{PNJ Importants}
\begin{center}
	\begin{tabular}{|P{2cm}|P{3cm}|P{3.5cm}|P{2.5cm}|}
	\hline
	\textbf{Nom} 			& \textbf{Rôle} 	& \textbf{Description} 								& \textbf{Motivation} \\ \hline
	Comte Orlov 			& Noble russe 		& 55 ans, élégant, froid. Porte un sceau impérial.	& Nostalgie : Veut restaurer la Russie impériale. \\ \hline
	Dr. Anastasia Romanova 	& Historienne 		& 30 ans, passionnée. A remarqué les anomalies. 	& Devoir : Veut protéger l'héritage russe. \\ \hline
	Inspecteur Petrov 		& Police 			& 45 ans, sceptique. Pense à un vol. 				& Devoir : Veut résoudre l'affaire. \\ \hline
	\end{tabular}
\end{center}

\subsubsection*{Menaces}
\begin{center}
	\begin{tabular}{|P{2cm}|P{3cm}|P{3.5cm}|P{2.5cm}|}
	\hline
	\textbf{Type} 			& \textbf{Description} 				& \textbf{Statistiques (DG)} 				& \textbf{Faiblesses} \\ \hline
	Soldats Fantômes 		& Soldats de l'époque impériale. 	& FOR 80\%, DEX 60\%, CON 70\% 				& Armes modernes, symboles soviétiques. \\ \hline
	Tsar Maudite (Esprit) 	& Peut geler les vivants. 			& FOR 90\%, DEX 80\%, CON 100\%, Magie 90\% & Feu, symboles communistes. \\ \hline
	Comte Orlov (corrompu) 	& Utilise des pouvoirs impériaux. 	& FOR 60\%, DEX 70\%, CON 60\%, Magie 80\% 	& Détruire le sceau impérial. \\ \hline
	\end{tabular}
\end{center}

\newpage

\section{Tableau Récapitulatif des Scénarios}
\begin{center}
\begin{small}
\rotatebox{90}{
	\begin{tabular}{|P{2.5cm}|P{2.0cm}|P{2.0cm}|P{3cm}|P{2.0cm}|P{2.0cm}|}
	\hline
	\textbf{Musée} 					& \textbf{Titre} 						& \textbf{Niveau} 			& \textbf{Menace Principale} 			& \textbf{Lieu Clé} 		& \textbf{Récompense} \\ \hline
	MNHN (Paris) 					& \textit{L'Éveil des Oubliés} 			& Intermédiaire 			& Os Vivants, Ombre de X-47 			& Cabinet des Curiosités 	& \textit{Codex Bestiarum} \\ \hline
	Louvre (Paris) 					& \textit{Le Souffle de la Joconde} 	& Débutant-Intermédiaire 	& Femme Sans Yeux, Statues Animées 		& Salle des Caryatides 		& Vraie Joconde (talisman) \\ \hline
	Musée de l'Homme (Paris) 		& \textit{Les Masques de la Folie} 		& Avancé 					& Démon des Masques, Masques Animés 	& Galerie des Crânes 		& Masque Dogon \\ \hline
	Quai Branly (Paris) 			& \textit{Les Tambours de la Mort} 		& Avancé 					& Esprits Ancestraux, Tambours Maudits 	& Réserves des Instruments 	& Gri-gri \\ \hline
	Pergamonmuseum (Berlin) 		& \textit{La Porte des Dieux Oubliés} 	& Avancé 					& Statues Animées, Araignée de Tiamat 	& Porte d'Ishtar 			& Parchemin en cunéiforme \\ \hline
	Neues Museum (Berlin) 			& \textit{Le Buste qui Chuchote} 		& Intermédiaire 			& Néfertiti (Démon), Momies Vivantes 	& Salle de Néfertiti 		& Scarabée du Nil \\ \hline
	DDR Museum (Berlin) 			& \textit{Les Archives de l'Oubli} 		& Débutant-Intermédiaire 	& Espions Fantômes, Ombre Vivante 		& Archives secrètes 		& Dossier secret de la Stasi \\ \hline
	British Museum (Londres) 		& \textit{Le Sceau de Salomon} 			& Avancé 					& Asmodée, Golems de Chair 				& Salle des Trésors 		& Sceau de Salomon \\ \hline
	Prado (Madrid) 					& \textit{Le Tableau qui Regarde} 		& Intermédiaire 			& Créatures Peintes, Diego López 		& Salle de Goya 			& Palette de Goya \\ \hline
	Rijksmuseum (Amsterdam) 		& \textit{Le Portrait du Diable} 		& Débutant 					& Démons Mineurs, Le Diable 			& Galerie d'Honneur 		& Portrait du Diable \\ \hline
	Ermitage (Saint-Pétersbourg) 	& \textit{La Malédiction des Tsars} 	& Avancé 					& Tsar Maudit, Soldats Fantômes 		& Salle des Trônes 			& OEuf de Fabergé maudit \\ \hline
	\end{tabular} }
\end{small}
\end{center}

\newpage

\section{Conseils pour le Maître de Jeu}
\subsection{Adaptation aux Règles de Delta Green / Félicie / EORA}
\begin{itemize}
\item \textbf{Compétences utiles} :
	\begin{itemize}
		\item \textbf{Occultisme} : Pour reconnaître les symboles, rituels, ou entités.
		\item \textbf{Histoire} : Pour comprendre le contexte historique des artefacts.
		\item \textbf{Langues} : Indispensable pour les scénarios à l'étranger (ex : allemand pour Berlin, russe pour Saint-Pétersbourg).
		\item \textbf{Psychologie} : Pour résister aux effets mentaux (perte de SAN, possession).
	\end{itemize}

\item \textbf{Équipement recommandé} :
	\begin{itemize}
		\item Armes bénites (balles en argent, lames en fer froid).
		\item Objets de protection (gri-gri, sceaux, talismans).
		\item Outils d'enquête (caméras thermiques, détecteurs de magie).
	\end{itemize}
\end{itemize}

\subsection{Liens entre les Scénarios}
Pour créer une \textbf{campagne cohérente}, tu peux lier ces scénarios entre eux :
\begin{itemize}
\item \textbf{Exemple 1 : \textit{La Conspiration des Musées}} :
	\begin{itemize}
		\item \textit{L'Éveil des Oubliés} (MNHN) $\rightarrow$ Les agents découvrent un lien avec le Louvre (\textit{Le Souffle de la Joconde}).
		\item \textit{Le Souffle de la Joconde} $\rightarrow$ Le Comte de Montclair a des connexions avec le British Museum (\textit{Le Sceau de Salomon}).
		\item \textit{Le Sceau de Salomon} $\rightarrow$ Le Rabin Eliézer travaille avec le culte de l'Araignée de Tiamat (Pergamonmuseum).
	\end{itemize}

\item \textbf{Exemple 2 : \textit{Le Réseau des Collectionneurs Maudits}} :
	\begin{itemize}
		\item \textit{Le Portrait du Diable} (Rijksmuseum) $\rightarrow$ Le Comte van der Meer a des liens avec le Comte Orlov (\textit{La Malédiction des Tsars}).
		\item \textit{La Malédiction des Tsars} $\rightarrow$ Le Comte Orlov a acheté des artefacts au Musée de l'Homme (\textit{Les Masques de la Folie}).
	\end{itemize}
\end{itemize}

\subsection{Idées pour des Campagnes Longues}
\begin{itemize}
\item \textbf{La Piste des Artefacts Maudits} :
	Les joueurs doivent retrouver \textbf{5 artefacts} (un dans chaque musée) avant qu'un culte ne les utilise pour réveiller un Grand Ancien. \emph{Musées concernés : Louvre, British Museum, Pergamonmuseum, Ermitage, Quai Branly. }

\item \textbf{La Guerre des Agences Occultes} :
	Félicie, EORA, et Delta Green se disputent le contrôle d'un artefact maudit (ex : le \textit{Sceau de Salomon}).
	Les joueurs doivent choisir un camp\ldots~  ou travailler en secret pour tous.

\item \textbf{Le Projet Muséum} :
	Un milliardaire fou (PNJ : Richard Van Houten) veut créer un \textbf{musée des horreurs} en collectionnant des artefacts maudits du monde entier.
	Les joueurs doivent l'arrêter avant qu'il n'ouvre le musée (et ne libère toutes les entités).
\end{itemize}


\subsection{Conseils pour l'Immersion}
\begin{itemize}
\item \textbf{Ambiance sonore} :
	\begin{itemize}
		\item Bruits de pas dans les musées vides.
		\item Chuchotements dans une langue ancienne (latin, égyptien, akkadien).
		\item Musique : Utilise des bandes-son de musées hantés (ex : \textit{The Haunted Museum} sur YouTube).
	\end{itemize}

\item \textbf{Ambiance visuelle} :
	\begin{itemize}
		\item Lumière : Les ampoules clignotent près des artefacts maudits.
		\item Reflets : Les vitrines montrent des images différentes selon l'angle.
		\item Ombres : Les ombres bougent toutes seules.
	\end{itemize}

\item \textbf{Effets spéciaux} :
	\begin{itemize}
		\item Fausses preuves : Donne aux joueurs des photos floues ou des enregistrements audio avec des voix étranges.
		\item Objets maudits : Utilise de vrais artefacts (ex : une pièce de monnaie ancienne) pour immerger les joueurs.
	\end{itemize}
\end{itemize}

%  Annexes 
\chapter*{Annexes}
\addcontentsline{toc}{chapter}{Annexes}

\section{Glossaire}
\begin{itemize}
	\item \textbf{Ankou} : Personnification de la Mort dans le folklore breton.
	\item \textbf{Byakhee} : Créature lovecraftienne, sorte de cheval démoniaque.
	\item \textbf{Chaugnar Faugn} : Entité lovecraftienne, souvent associée à des cultes déviants.
	\item \textbf{Contrées du Rêve} : Dimension parallèle où règne le chaos.
	\item \textbf{DGSE} : Direction Générale de la Sécurité Extérieure (service de renseignement français).
	\item \textbf{EORA} : European Occult Response Agency (agence européenne de lutte contre l'occulte).
	\item \textbf{Félicie} : Agence secrète française de lutte contre l'occulte.
	\item \textbf{Mi-Go} : Créatures extraterrestres de l'univers lovecraftien.
	\item \textbf{Nightgaunt} : Créature des Contrées du Rêve, souvent décrite comme une ombre ailée.
	\item \textbf{Shoggoth} : Entité amorphe et maléfique de l'univers lovecraftien.
\end{itemize}

\newpage

%  Notes 
\section{Notes}

Ce livret est conçu pour être utilisé avec les règles de \textbf{Delta Green}. Les descriptions, scénarios, et fiches de personnages sont inspirés des travaux de \textbf{Tristan Lhomme} (notamment l'agence Félicie dans \textit{Le Musée de Lhomme}) et de l'univers de \textbf{Delta Green}.

N'hésitez pas à \textbf{adapter et personnaliser} le contenu pour vos propres parties. Les idées de scénarios, les fiches de personnages, et les descriptions de lieux sont là pour vous inspirer et vous aider à créer des aventures mémorables.

\begin{shaded}
	\begin{center}
		\textbf{Bonne partie, et que les Grands Anciens vous épargnent !}
	\end{center}
\end{shaded}

%% \newpage

%  Sources et Références 
% \chapter*{Sources et Références}
% \addcontentsline{toc}{chapter}{Sources et Références}

\addcontentsline{toc}{section}{Bibliographie}

\begin{thebibliography}{9}

	\bibitem{legrog}
	Le Grog. (s.d.). \textit{Musée de Lhomme}. \url{https://www.legrog.org/jeux/appel-de-cthulhu/musee-de-lhomme-fr}
	
	\bibitem{ludovox}
	LudoVox. (2015). \textit{[JDR] L'appel de Cthulhu : Le musée de l'homme}. \url{https://ludovox.fr/jdr-lappel-de-cthulhu-le-musee-de-lhomme/}
	
	\bibitem{unification}
	Unification France. (2015). \textit{Jeu de rôle l'Appel de Cthulhu : Le Musée de Lhomme}. \url{https://www.unificationfrance.com/article39226.html}
	
	\bibitem{agorajeux}
	Agorajeux. (s.d.). \textit{Le Musée de Lhomme}. \url{https://www.agorajeux.com/fr/cthulhu/5349-le-musee-de-lhomme.html}
	
	\bibitem{sansdetour}
	Sans-Détour éditions. (2014). \textit{Le Musée de Lhomme}. \url{https://www.sans-detour.net/}
	
	\bibitem{delta-green}
	Arc Dream Publishing. (2016). \textit{Delta Green: Agent's Handbook}. 
	
	\bibitem{cthulhu}
	Chaosium. (1981). \textit{L'Appel de Cthulhu} (6ème édition).

\end{thebibliography}

%% \newpage

\end{document}
